% !TEX program = xelatex
\documentclass[12pt,twoside]{report} %openright
	%\usepackage[utf8]{inputenc}
 	\usepackage[greek,english]{babel}
    \usepackage{yfonts}
	\usepackage{latexsym}
	\usepackage{amsmath}
	\usepackage{amssymb}
	\usepackage{amsthm}
	\usepackage{stmaryrd}
	\usepackage{pb-diagram}
	\usepackage{amscd}
	\usepackage{mathrsfs}
	\usepackage{enumerate}
	\usepackage{color}
	\usepackage{cite}
	%\usepackage{graphicx}
	\usepackage[all]{xy}
	
	 
	\usepackage{graphicx}
\graphicspath{ {./images/} }
	
	%\usepackage{fixltx2e}
	\usepackage{extarrows}
	\usepackage{datetime}
	\usepackage{mathtools}
    \usepackage{remreset}
    \usepackage{afterpage}
    \usepackage{verbatim}

    \usepackage[svgnames]{xcolor}
	\usepackage{hyperref}
\hypersetup{
    colorlinks=true,
    linkcolor=[RGB]{61,82,62},
    filecolor=[RGB]{61,82,62},
    urlcolor=[RGB]{61,82,62},
    citecolor=[RGB]{61,82,62},
    pdftitle={Groupoïde fondamental et théorème de Van Kampen en théorie des topos},
    pdfauthor={O. Leroy},
    pdflang={fr}
}

\usepackage{bookmark}
    %\usepackage{csquotes}
    \usepackage{etoc}

\usepackage{pdflscape}
	
	\usepackage{fontspec}
	\setmainfont{EB Garamond}
	\usepackage{float}
	\usepackage[12pt]{moresize}
	\usepackage{titlesec}
	\usepackage[twoside,top=1.5in,bottom=1.5in,left=1in,right=1.3in,headheight=16pt,headsep=30pt,footskip=40pt]{geometry}

	\linespread{1.25}
%	\linespread{1.05}


%\NeedsTeXFormat{LaTeX2e}
%\ProvidesPackage{quiver}[2021/01/11 quiver]
\RequirePackage{tikz-cd}
\RequirePackage{amssymb}
\usetikzlibrary{calc}
\usetikzlibrary{decorations.pathmorphing}


% A TikZ style for curved arrows of a fixed height, due to AndréC.
\tikzset{curve/.style={settings={#1},to path={(\tikztostart)
    .. controls ($(\tikztostart)!\pv{pos}!(\tikztotarget)!\pv{height}!270:(\tikztotarget)$)
    and ($(\tikztostart)!1-\pv{pos}!(\tikztotarget)!\pv{height}!270:(\tikztotarget)$)
    .. (\tikztotarget)\tikztonodes}},
    settings/.code={\tikzset{quiver/.cd,#1}
        \def\pv##1{\pgfkeysvalueof{/tikz/quiver/##1}}},
    quiver/.cd,pos/.initial=0.35,height/.initial=0}

% TikZ arrowhead/tail styles.
\tikzset{tail reversed/.code={\pgfsetarrowsstart{tikzcd to}}}
\tikzset{2tail/.code={\pgfsetarrowsstart{Implies[reversed]}}}
\tikzset{2tail reversed/.code={\pgfsetarrowsstart{Implies}}}
% TikZ arrow styles.
\tikzset{no body/.style={/tikz/dash pattern=on 0 off 1mm}}

	% use osf for figures in math mode
	\DeclareSymbolFont{digits}{\encodingdefault}{\rmdefault}{m}{n}
	\SetSymbolFont{digits}{normal}{\encodingdefault}{\rmdefault}{m}{n}
	\DeclareMathSymbol{0}{\mathalpha}{digits}{"30}
	\DeclareMathSymbol{1}{\mathalpha}{digits}{"31}
	\DeclareMathSymbol{2}{\mathalpha}{digits}{"32}
	\DeclareMathSymbol{3}{\mathalpha}{digits}{"33}
	\DeclareMathSymbol{4}{\mathalpha}{digits}{"34}
	\DeclareMathSymbol{5}{\mathalpha}{digits}{"35}
	\DeclareMathSymbol{6}{\mathalpha}{digits}{"36}
	\DeclareMathSymbol{7}{\mathalpha}{digits}{"37}
	\DeclareMathSymbol{8}{\mathalpha}{digits}{"38}
	\DeclareMathSymbol{9}{\mathalpha}{digits}{"39}
	% helvetica for sans in math
	\DeclareMathAlphabet{\mathsf}{\encodingdefault}{phv}{m}{n}
	%\usepackage{showkeys}
	%\usepackage{showframe}
	\usepackage{fancyhdr}
	%\fancyhead[CO]{\rightmark}
	\fancyhead[CO]{}
	%\fancyhead[CE]{\leftmark}
	\fancyhead[CE]{}
	\fancyhead[RE]{}
	\fancyhead[LO]{}
	\fancyhead[RO]{}
	\fancyhead[LE]{}
	\cfoot{\thepage}
	\renewcommand{\headrulewidth}{0.0pt}
	\titleformat\section{\normalfont\Large\centerline{\hss\rule{1.3in}{.01in}\hss}}{\thesection}{1em}{}
    \titleformat{\chapter}[display]{}{\huge\filcenter { \thechapter}}{12pt}{\large \filcenter}

%\newcommand\hmmax{0}
%\newcommand\bmmax{0}

%%%%%%%%%%%%%%%%%%%%%%%%%%%%%%%%%%%%%%%%%%%%%%


\newcommand*\isommap{%
  \xrightarrow{\raisebox{-0.2 em}{\smash{\ensuremath{\sim}}}}%
}

\newcommand*\isomlong{%
  \xlongrightarrow{\raisebox{-0.2 em}{\smash{\ensuremath{\sim}}}}%
}

 \newcommand{\cHom}{\mathrel{\mathscr{H}\mkern-2mu\text{\upshape om}}}
 \newcommand{\cTor}{\mathrel{\mathscr{T}\mkern-2mu\text{\upshape or}}}
 \newcommand{\cExt}{\mathrel{\mathscr{E}\mkern-2mu\text{\upshape xt}}}

   \DeclareMathOperator{\D}{D}
    \DeclareMathOperator{\K}{K}
    \DeclareMathOperator{\Ld}{\mathrel{L}}
    \DeclareMathOperator{\Rd}{\mathrel{R}}
    \DeclareMathOperator{\Hom}{Hom}
    \DeclareMathOperator*{\bHom}{\mathnormal{\mathbf{Hom}}}
    \DeclareMathOperator*{\bRd}{\mathrm{\mathbf{R}}} 
    \DeclareMathOperator*{\bLd}{\mathrm{\mathbf{L}}}
    \DeclareMathOperator{\bH}{\mathbf{H}}
    \DeclareMathOperator{\Fibtop}{Fibtop}
    \DeclareMathOperator{\SLC}{SLC}
    \DeclareMathOperator{\Ob}{Ob}
    \DeclareMathOperator{\Cart}{Cart}
    \DeclareMathOperator{\Cartop}{Cartop}
    \DeclareMathOperator{\Top}{Top}
    \DeclareMathOperator{\Grpd}{Grpd}

    \DeclareMathOperator{\Ens}{Ens}
    \DeclareMathOperator{\Dis}{Dis}
    \DeclareMathOperator{\Aut}{Aut}
    \DeclareMathOperator{\Homtop}{Homtop}
    \DeclareMathOperator{\pr}{pr}
    \DeclareMathOperator{\Point}{Point}
    \DeclareMathOperator{\SLCF}{SLCF}
    \DeclareMathOperator{\LC}{LC}
    \DeclareMathOperator{\App}{App}
    \DeclareMathOperator{\Img}{Img}
    \DeclareMathOperator{\Fonct}{Fonct}
    \DeclareMathOperator{\cart}{cart}
    \DeclareMathOperator{\Tors}{Tors}
    \DeclareMathOperator{\Isom}{Isom}
    \DeclareMathOperator{\Card}{Card}


	\def\cU{{\cal U}}
	\def\cG{{\cal G}}
	\newcommand{\isom}{\simeq}

%%%%%%%%%%%%%%%%%%%%%%%%%%%%%%%%%%%%%%%%%%%%%%

\makeatletter
\newcommand*{\toccontents}{\@starttoc{toc}}
\makeatother

\makeatletter
\@addtoreset{footnote}{section}
\makeatother

%\usepackage{sectsty}
%\partfont{\LARGE}

\begin{document}
\setcounter{tocdepth}{2}
\sloppy

\newgeometry{left=1.5in,right=1.5in, bottom=1.5in}

\thispagestyle{empty}
%\vspace*{\fill}
\noindent\makebox[\linewidth]{\rule{4in}{0.01in}}
\vspace*{\fill}
\begin{center}
    {\Huge Groupoïde Fondamental et} \\
    \vspace{0.5em}
    {\Huge Théorème de Van Kampen} \\
    \vspace{0.5em}
    {\Huge en Théorie des Topos} \\
    \vspace{0.5em}
    {\Large par} \\
    \vspace{0.5em}
    {\huge Olivier Leroy}
\end{center}
\vspace*{\fill}
\noindent\makebox[\linewidth]{\rule{4in}{0.01in}}
%\vspace*{\fill}

\newpage 
\thispagestyle{empty}

\null\vfill
\noindent\begin{minipage}{\textwidth}
This project, involving the transcription of J.-P. Jouanolou's thesis, was conducted with his authorization. It was carried out under the direction of Mateo Carmona, with the collaboration of Niels Borne and other volunteers. The transcription aims to be as faithful as possible to the original. This edition should be regarded as provisional; remarks, comments, and corrections are welcome. More details are available at: \url{https://github.com/gmateocarmona/categories-derivees-en-cohomologie-l-adique}
\end{minipage}

\newpage

\renewcommand{\contentsname}{TABLE DES MATIÈRES}
\tableofcontents\thispagestyle{empty}

\newpage

\thispagestyle{empty}
\begin{center}
\setlength{\baselineskip}{24pt}
\  \\
{
\ \\
[0.5in] CAHIERS MATHÉMATIQUES MONTPELLIER \\ 1979 \\
[1.5in] {\large \bf GROUPOÏDE FONDAMENTAL ET THÉORÈME DE VAN KAMPEN EN THÉORIE DES TOPOS} \\ Olivier LEROY \\
[2in] UNIVERSITÉ DES SCIENCES \\ ET TECHNIQUES \\ DU LANGUEDOC \\ U. E. R. DE MATHÉMATIQUES \\
Place Eugène Bataillon \\ 34060 MONTPELLIER CEDEX 
}
\end{center}
\vfill
\vfill
\vfill

\newpage

\begin{center}
\thispagestyle{empty}
\setlength{\baselineskip}{24pt}
\  \\
{\large
\ \\
[3in] Ce travail m'a été suggéré par \emph{C. Contou-Carrère} et \emph{A. Grothendieck}. \\
Mes résultats ont été exposés lors de séances de géométrie algébrique faites avec \emph{D. Alibert} et \emph{C. Contou-Carrère}. \\
}
\vfill
\end{center}
\vfill

\newpage

%%%%%%%%%%%%%%%%%%%%%%%%%%%%%%%%%%%%%%%%%%%%%%%%%%%%%%%%%%%%%%%
\chapter*{TABLE ANALYTIQUE}\thispagestyle{empty}
\addcontentsline{toc}{chapter}{Table analytique}
\section*{}

\begin{itemize}
    \item[1.] Objets connexes dans un topos
    \item[2.]
    \item[3.]
    \item[4.]
    \item[A.] {\bf Appendice}
    \begin{itemize}
        \item[ ] Catégories fibrées en topos 
    \end{itemize}
    \item[5.]
    \begin{itemize}
        \item[5.1.] 
        \item[5.2.] 
    \end{itemize}
    \item[ ] {\bf Références}
\end{itemize}

\newpage
\thispagestyle{empty}
\mbox{}

%%%%%%%%%%%%%%%%%%%%%%%%%%%%%%%%%%%%%%%%%
\setcounter{page}{0}

% Begin


%%%%%%%%%%%%%%%%%%%%%%%%%%%%%%%%%%%%%%%%%%%%%%%%%%%%%%%%%%%%%%%
\chapter*{\S \space I. --- OBJETS CONNEXES DANS UN TOPOS}
\addcontentsline{toc}{chapter}{I. Objets connexes dans un topos}
\section*{}

Tous les topos considérés sont des $\cU$-topos.
\vskip .3cm
{
Définitions {\bf 1.1.} --- \it
\begin{enumerate}
    \item[a)] Un objet d'un topos est \emph{connexe} s'il n'est pas somme directe de deux objets non-vides.
    \item[b)] Soit $X$ un objet d'un topos. On appelle \emph{composante connexe} de $X$ tout sous-objet connexe et non-vide $C$ de $X$ tel que $X$ soit somme directe de $C$ et d'un autre objet.
    \item[c)] Un \emph{topos} est connexe si son objet final est connexe.
    \item[d)] Un topos est \emph{localement connexe} s'il est engendré par ses objets connexes. 
\end{enumerate}
}
\vskip .3cm
{
Proposition {\bf 1.2.} --- \it Soit $C$ un objet d'un topos $E$. Les propriétés suivantes sont équivalentes~:
\begin{enumerate}
    \item[a)] $C$ est connexe et non-vide.
    \item[b)] Le foncteur
    $$
    \Hom_E(C, -): E \to \Ens 
    $$
    commute aux sommes directes.
    \item[c)] Pour tout ensemble $I$, l'application naturelle $I \to \Hom(C, I_E)$ est bijective (c'est immédiat).
\end{enumerate}
}
\vskip .3cm
{
Proposition {\bf 1.3.} --- \it Soit $(U_i \xlongrightarrow{f_i} V)_{i \in I}$ une famille épimorphique d'un topos $E$. 

Considérons les propriétés~:
\begin{enumerate}
    \item[(a)] $V$ est connexe et non-vide.
    \item[(b)] Le graphe $R \subset I \times I$ de la relation
    $$
    ``U_i \times_V U_j \quad \text{n'est pas vide''}
    $$
    est connexe (en tant que graphe ayant $I$ pour ensemble de sommets). On a
\end{enumerate}   
\begin{enumerate}
    \item[(i)] si les $U_i$ sont non-vides, (a) entraîne (b).
    \item[(ii)] si les $U_i$ sont connexes, et non-vides, (b) entraîne (a).
\end{enumerate}
}
\vskip .3cm
Démonstration. C'est trivial si $I = \emptyset$. On suppose donc $I \neq \emptyset$.
\vskip .3cm
{\bf $(a) \Rightarrow (b)$} ($U_i$ non-vides). Soit $(I_1, I_2)$ une partition de $I$ telle que pour tout $i \in I_1$ et tout $j \in I_2$, $U_i \times_V U_j$ soit vide. Si on désigne par $V_1$ et $V_2$ respectivement les images des morphismes
$$
\amalg_{i \in I_1} U_i \to V, \quad \amalg_{i \in I_2} U_i \to V
$$
alors $V$ est somme de $V_1$ et $V_2$. Donc $V_1$ ou $V_2$ est vide. Comme aucun
$U_i$ n'est vide, $I_1$ ou $I_2$ est vide.
\vskip .3cm
{\bf $(b) \Rightarrow (a)$} ($U_i$ connexes et non-vides). Soit $(Y_{\alpha})_{\alpha \in A}$ une famille d'objets de $E$, et considérons un morphisme
$$
V \to Y = \amalg_{\alpha} Y_{\alpha}.
$$
Pour chaque $\alpha \in A$, soit $I_{\alpha}$ l'ensemble des $i \in I$ tels que le composé
$$
U_i \to V \to Y
$$
se factorise par $Y_{\alpha}$. Puisque les $U_i$ sont connexes et non-vides, $I$ est réunion disjointe des $I_{\alpha}$. Soient $\alpha$ et $\beta$ deux indices distincts. Si $i \in I_{\alpha}$ et $j \in I_{\beta}$, $U_i \times_V U_j$ est vide puisque c'est un sous-objet de $U_i \times_Y U_j$. Appliquant (b), on voit que $I = I_{\alpha_0}$ pour un $\alpha_0 \in A$ et un seul. Donc $V$ est connexe et non-vide par $(1.2).$
\vskip .3cm
{
Proposition {\bf 1.4.} --- \it Tout objet d'un topos localement connexe est somme directe d'objets connexes (donc somme directe de ses composantes connexes).
}
\vskip .3cm
Démonstration. Soit $E$ un topos localement connexe. Soient $Y$ un objet de $E$ et $(U_i \to Y)_{i \in I}$ une famille épimorphique de $E$, où les $U_i$ sont connexes et non-vides. Soit $R$ le graphe de la relation $``U_i \times U_j$ n'est pas vide''. Pour chaque composante connexe $r$ de $R$, soit $C_r$ l'image dans $Y$ de la somme des $U_i$, $i$ parcourant l'ensemble des $i \in I$ qui sont sommets de $r$. $Y$ est somme directe des $C_Y$, qui sont connexes et non-vides d'après (1.3).
\vskip .3cm
{
Proposition {\bf 1.5.} --- \it Pour qu'un topos $E$ soit localement connexe, il faut et il suffit que le foncteur
\begin{align*}
  \Ens \to E \\
  I \mapsto I_E
\end{align*}
admette un adjoint à gauche
$$
c: E \to \Ens.
$$
Dans ce cas, étant donné un objet $X$ de $E$, les produits fibrés
\[
\xymatrix{
	X_\gamma \ar[r] \ar[d]  & e_E \ar[d]^{\gamma} \\
	X\ar[r]                       & c(X)_E
}
\]
($\gamma$ parcourant $c(X)$) sont les composantes de $X$.
}
\vskip .3cm
Démonstration.
\begin{enumerate}
    \item[$(i)$] Supposons $E$ localement connexe. Pour tout objet $X$ de $E$, désignons par $c(X)$ l'ensemble des classes de $X$-isomorphisme de composantes connexes de $X$ (cet ensemble est bien sûr $\cU$-petit). Soit $f: X \to Y$ un morphisme de $E$~; étant donnée une composante connexe $C$ de $X$, il existe une composante connexe $D$ de $Y$, unique à $Y$-isomorphisme près, telle que $f_{/C}$ se factorise par $D$. D'où une application
    $$
    c(X) \to c(Y).
    $$
    On a ainsi obtenu un foncteur covariant
    $$
    c: E \to \Ens.
    $$
    Le foncteur $c$ est adjoint à gauche de $I \mapsto I_E$: en effet, étant donnés un objet $X$ de $E$ et un ensemble $I$, on définit une application~: 
    $$
    \App(c(X), I) \to \Hom(X, I_E)
    $$
    en associant à l'application
    $$
    a: c(X) \to I
    $$
    le morphisme $X \to I_E$ dont la restriction à chaque composante connexe $C$ de $X$ est la section de $I_E$ au-dessus de $C$ définie par $a(C) \in I$~; cette application est bijective par (1.2), et elle est fonctorielle en $X$ et $I$.
    \item[$(ii)$] Inversement, supposons qu'on ait un adjoint à gauche $c: E \to \Ens$ du foncteur $I \to I_E$.
    
    Soit $X$ un objet de $E$. Avec les notations de l'énoncé, $X$ est somme directe des $X_{\gamma}$, $\gamma \in c(X)$. Il suffit donc de prouver que les $X_{\gamma}$ sont connexes et non-vides. Or, pour tout ensemble $I$, les applications naturelles
    $$
    I \to \Hom(X_{\gamma}, I_E)
    $$
    fournissent une application
    $$
    I^{c(X)} \to \prod_\gamma \Hom(X_{\gamma}, I_E)
    $$
    qui rend commutatif le diagramme
    \[
        \xymatrix{
        I^{c(X)} \ar[d] \ar[dr]^-\sim \\
        \prod_\gamma \Hom(X_{\gamma}, I_E) \ar[r]^-\sim & \Hom(X, I_E)
        }
    \]
    donc chacune des applications
    $$
    I \to \Hom(X_{\gamma}, I_E)
    $$
    est bijective~; on conclut par (1.2).
\end{enumerate}


% End

% Begin


%%%%%%%%%%%%%%%%%%%%%%%%%%%%%%%%%%%%%%%%%%%%%%%%%%%%%%%%%%%%%%%
\chapter*{\S \space II. --- OBJETS LOCALEMENT CONSTANTS ET OBJETS GALOISIENS}
\addcontentsline{toc}{chapter}{II. Objets localement constants et objets galoisiens}
\section*{}

On se donne un $\cU$-topos $E$.

{\bf 2.1. Objets localement constants}
\vskip .3cm
{
Définitions {\bf 2.1.1.} --- \it
\begin{enumerate}
    \item[1)] Soient $A$ et $B$ deux objets de $E$. On dit que $A$ trivialise $B$ si $A \times B \xlongrightarrow{pr_1} A$ est un objet constant du topos $E_{/A}$.
    \item[2)] Nous dirons qu'un objet $L$ de $E$ est \emph{localement constant} si les objets de $E$ qui trivialisent $L$ recouvrent $E$.
    \item[3)] Enfin nous dirons qu'un préfaisceau $F$ sur une catégorie $C$ est localement constant si pour tout morphisme $X \to Y$ de $C$, l'application $F(Y) \to F(X)$ est bijective.
\end{enumerate}
}
\vskip .3cm
{
Proposition {\bf 2.1.2.} --- \it Soient $Y$ un objet de $E$ et $D \xlongrightarrow{u} C$ un morphisme de $E$. 
\begin{enumerate}
    \item[(i)] Si $C$ trivialise $Y$, $D$ trivialise $Y$ ;
    \item[(ii)] si, en outre, $C$ et $D$ sont connexes et non-vides, l'application $f \mapsto f \circ u$ de $\Hom(C, Y)$ dans $\Hom(D, Y)$ est \emph{bijective}.
\end{enumerate}
}
\vskip .3cm
Démonstration. --- Prenons un ensemble $I$ et un $C$-isomorphisme $\alpha: C \times Y \isom I_C$.
\begin{enumerate}
    \item[(i)] On a des diagrammes cartésiens
    $$
        \xymatrix{
            D\times Y \ar[r] \ar[d] & C\times Y \ar[d]\\
            D \ar[r] & C
        }
        \qquad \qquad 
        \xymatrix{
            I_D \ar[r] \ar[d] & I_C \ar[d]\\
            D \ar[r] & C
        }
    $$
    d'où un $D$-isomorphisme
    $$
    \beta: D \times Y \to I_D.
    $$
    \item[(ii)] Soient $f: C \to Y$ et $g: D \to Y$. Pour que $f \circ u = g$, il faut et il suffit que le diagramme
    $$
    \xymatrix{
        D\times Y \ar[r] & C\times Y\\
        D \ar[u]^{(1_D,g)} \ar[r] & C \ar[u]_{(1_C,f)}
    }
    $$
    soit commutatif, ou encore que soit commutatif le diagramme
    $$
        \xymatrix{
            I_D \ar[r] & I_C \\
            D \ar[r] \ar[u]^{\beta\circ (1_D,g)}
            & C \ar[u]_{\alpha\circ (1_C,f)}
        }
    $$
    d'où le point $(ii)$ par (1.3).
\end{enumerate}

\vskip .3cm
{
Proposition {\bf 2.1.3.} --- \it 
\begin{enumerate}
    \item[(i)] Soient $C$ et $X$ deux objets de $E$. On a un morphisme naturel
    $$
    p: \Hom(C, X)_C \to X
    $$
    c'est le morphisme qui, pour tout $f \in \Hom(C, X)$, rend commutatif le diagramme
    $$
        \xymatrix{
            C \ar[d]_{i_f} \ar[dr]^f \\
            \Hom(C, X)_C \ar[r]_-p & X
        }
    $$
    où $i_f$ désigne la section de $\Hom(C, X)_C$ au-dessus de $C$ définie par $f$.

    On a donc aussi un $C$-morphisme naturel :
    $$
    m = (q, p): \Hom(C, X)_C \to C \times X
    $$
    où $q$ désigne le $C$-morphisme $\Hom(C, X)_C \to C$.
    \item[(ii)] Si $C$ est connexe non-vide et trivialise $X$, le $C$-morphisme naturel
    $$
    \Hom(C, X)_C \to C \times X
    $$
    est un \emph{isomorphisme}.
\end{enumerate}
}
\vskip .3cm

{\it Démonstration}. En effet, prenons un ensemble $I$ et un $C$-isomorphisme $I_C \to C \times X$. On en tire une bijection :
$$
\Hom_C(C, I_C) \isom \Hom_C(C, C \times X) \isom \Hom(C, X)
$$
d'où un diagramme de $E_{/C}$~:
$$
    \xymatrix{
        (\Hom_C(C, I_C))_C \ar[r]^-\sim \ar[d] & \Hom(C, X)_C \ar[d]\\
        I_C \ar[r] & C\times X
    }
$$
dont la commutativité prouve notre assertion.

\vskip .3cm
{
Proposition {\bf 2.1.4.} --- \it Soient $p$ un point de $E$, $C$ un objet connexe de $E$, $y_0$ un point de la fibre $p^{-1}(C)$ et $X$ un objet de $E$. Si $C$ trivialise $X$, l'application
\begin{align*}
\Hom(C, X) & \to p^{-1}(X)\\
f & \mapsto f(y_0)
\end{align*}
est \emph{bijective}.
}
\vskip .3cm

{\it Démonstration}. En effet, elle se déduit de l'application
\begin{align*}
p^{-1}(C) \times\Hom(C, X) & \to p^{-1}(C) \times p^{-1}(X))\\
(y, f) & \mapsto (y, f(y))
\end{align*}
qui provient, par passage aux fibres, de l'isomorphisme naturel
$$
\Hom(C, X)_C \to C \times X \quad (2.1.3, ii)
$$

\vskip .3cm
{
Proposition {\bf 2.1.5.} --- \it Soient $L$ un objet localement constant de $E$ et $U$ l'image de $L \to e_E$. Il existe un objet $V$ de $E$ tel que $U \amalg V$ soit isomorphe à $e_E$.
}
\vskip .3cm

{\it Démonstration}. Recouvrons $e_E$ par des objets $(U_i)_{i \in I}$ qui trivialisent $L$. Soit $I_0$ (resp. $I_1$) l'ensemble des $i \in I$ tels que $U_i \times L$ soit vide (resp. non-vide). Pour tout $i \in I_0$ et tout $j \in I_1$, $U_i \times U_j$ est vide; en effet, il existe un ensemble non-vide $F$ tel que $F_{U_i \times U_j}$ soit vide.

Soient $V_0$ et $V_1$ respectivement les images de :
$$
\amalg_{i \in I_0} U_i \to e_E, \quad \amalg_{i \in I_1} U_i \to e_E.
$$
Évidemment, $V_1$ est un sous-objet de $U$, et $U \times V_0$ est vide. Comme $e_E$ est somme de $V_0$ et $V_1$, on conclut que $V_1 \isom U$ et $U \amalg V_0 \isom e_E$. 

\vskip .3cm
{\bf 2.2.} Nous supposons maintenant le topos $E$ \emph{localement connexe}

\vskip .3cm
{
Proposition {\bf 2.2.1.} --- \it Soient $Z$ un objet de $E$ et $S$ une sous-catégorie génératrice de $E$ dont les objets sont connexes et non-vides dans $E$. Les propriétés suivantes sont équivalentes :
\begin{enumerate}
    \item[a)] Tout objet de $S$ trivialise $Z$ ;
    \item[b)] Pour tout morphisme $D \to C$ de $S$, l'application correspondante $\Hom(C, Z) \to \Hom(D, Z)$ est bijective ; autrement dit, $\Hom(-, Z)$ est un préfaisceau localement constant sur $S$.
\end{enumerate}
}
\vskip .3cm

{\it Démonstration}. (a) $\Rightarrow$ (b) : c'est (2.1.2). (b) $\Rightarrow$ (a) : soit $C_0$ un objet de $S$. Pour tout objet $D$ de $S$, l'application
\begin{align*}
\Hom(D, C_0) \times \Hom(C_0, Z) & \to \Hom(D, C_0) \times \Hom(D, Z)\\
(u, f) & \mapsto (u, f \circ u)
\end{align*}
au-dessus de $\Hom(D, C_0)$ est bijective, d'où un $C_0$-isomorphisme
$$
\Hom(C_0, Z)_{C_0} \to C_0 \times Z.
$$

\vskip .3cm
{
Proposition {\bf 2.2.2.} --- \it Soit $S$ une sous-catégorie génératrice de $E$ dont les objets sont connexes et non-vides dans $E$. Munissons $S$ de la topologie induite par $E$. Tout préfaisceau localement constant sur $S$ est un faisceau.
}
\vskip .3cm

{\it Démonstration}. Soit $F$ un préfaisceau localement constant sur $S$. Il est clair que $F$ est un préfaisceau séparé ; prouvons que c'est un faisceau.

Soient $U$ un objet de $S$ et $R$ un crible couvrant $U$ (i.e. $R$ contient une famille épimorphique de $E$) ; soit enfin une section
$$
t: R \to F.
$$
Pour tout objet $V$ de $S$ et toute $f: V \to R$, désignons par $x(V, f)$ la section de $F$ au-dessus de $U$, image inverse de $t(f) \in F(V)$ par la bijection $F(U) \to F(V)$ correspondant à $f: V \to U$. Étant donné un diagramme commutatif
$$
\xymatrix{
    V_0 \ar[d] \ar[dr]^{f_0}\\
    V_1 \ar[r]_{f_1} & R
}
$$
on a $x(V_0, f_0) = x(V_1, f_1)$ ; donc (d'après (1.3)) les $x(V, f)$ sont tous égaux à un même $x \in F(U)$ ; et on a $x_{/R} = t$ par construction de $x$.

\vskip .3cm
{
Définition auxiliaire {\bf 2.2.3.} --- \it
\begin{enumerate}
    \item[(i)] Soit $L$ un objet localement constant de $E$. Soit $S$ une sous-catégorie génératrice de $E$, dont les objets sont connexes non-vides dans $E$ et trivialisent $L$. On appellera \emph{site générateur adapté à} $L$ une telle sous-catégorie, munie de la topologie induite par $E$. 
    \item[(ii)] Si $(L_i)_{1 \leq i \leq n}$ est une famille finie d'objets localement constants de $E$, il existe un site générateur de $E$ adapté à chacun des $L_i$ : cela découle de (2.1.2, (i)) par récurrence sur $n$.
\end{enumerate}
}
\vskip .3cm

\vskip .3cm
{
Proposition {\bf 2.2.4.} ($\varprojlim$ et $\varinjlim$ d'objets localement constants). --- \it

Soient $I$ une petite catégorie et 
$L: I \to E$
un foncteur tel que $L(i)$ soit localement constant pour tout $i \in \Ob(I)$. \emph{Supposons qu'il existe un site générateur $S$ adapté à tous les $L(i)$} ; alors :
\begin{enumerate}
    \item[(i)] $P = \varprojlim\limits_{I} L(i)$ est localement constant et $S$ est un site générateur adapté à P (2.2.1).
    \item[(ii)] Le préfaisceau sur $S$ :
    $$
    C \to L(C) = \varinjlim_I \Hom (C, L(i))
    $$
    est un faisceau ; donc $\varinjlim\limits_{I} L(i)$ est un objet localement constant, isomorphe à $L$ en tant que faisceau sur $S$ (en particulier, $S$ est un site générateur adapté à $\varinjlim L(i)$).
    \item[(iii)] Corollaire : Soit $u: L \to M$ un morphisme entre objets localement constants de $E$, et soit $S$ un site générateur adapté à $L$ et $M$. En tant que faisceau sur $S$, l'image de $u$ est isomorphe à 
    $$
    C \to \Img(\Hom(C, L) \to \Hom(C, M)).
    $$
\end{enumerate}
}
\vskip .3cm

\vskip .3cm
{
Proposition {\bf 2.2.5.} --- \it Soit $L$ un objet localement constant de $E$.
\begin{enumerate}
    \item[(i)] Tout sous-objet localement constant $K$ de $L$ est somme directe de composantes connexes de $L$.
    \item[(ii)] Soient $C$ un objet connexe non-vide de $E$ et $K$ un sous-objet de $L$ qui est somme directe de composantes connexes de $L$. Si $C$ trivialise $L$, $C$ trivialise $K$.
\end{enumerate}
}
\vskip .3cm

{\it Démonstration}.
\begin{enumerate}
    \item[(i)] Soit $S$ un site générateur adapté à $K$ et à $L$. Le préfaisceau sur $S$
    $$
    C \to M(C) = \Hom(C, L)\ - \ \Hom(C, K)
    $$
    est un faisceau (2.2.2), et $L$ est somme directe de $K$ de $M$, c.q.f.d.
    \item[(ii)] Si $L$ est somme directe de deux sous-objets $K$ et $K'$, alors $C \times L$ qui est constant dans $E_{/C}$ est somme directe de $C \times K$ et $C \times K'$, qui sont donc constants dans $E_{/C}$ puisque $C$ est connexe et non-vide.
\end{enumerate}

\vskip .3cm
{
Corollaires {\bf 2.2.6.} --- \it 
\begin{enumerate}
    \item[(a)] Tout morphisme d'un objet localement constant non-vide de $E$ dans un objet localement constant connexe de $E$ est un épimorphisme.
    \item[(b)] Soient des morphismes de $E$ :
    $$
    \xymatrix{
        L \ar@<1ex>[r]^f \ar@<-1ex>[r]_g & M
    }
    $$
    où $L$ et $M$ sont localement constants et $L$ en outre, connexe. Si le noyau de $(f, g)$ est non-vide, alors $f = g$.
\end{enumerate}
}
\vskip .3cm

((a) découle de 2.2.4, (iii) et (b) de 2.2.4, (i)).

\vskip .3cm
{\bf 2.3. Objets galoisiens}

\emph{On ne suppose plus} le topos $E$ localement connexe.

\vskip .3cm
{\bf 2.3.1.} {Définition. --- \it Nous dirons qu'un objet $Y$ de $E$ est \emph{galoisien} s'il est localement constant, connexe et non-vide, et s'il est un pseudo-torseur sous le groupe constant $\Aut(Y)_E$.
}
\vskip .3cm
{\it Remarques}.
\begin{enumerate}
    \item[(i)] D'après (2.1.5.), cela revient à dire que l'image $U$ de $Y \to e_E$ est une composante connexe de $e_E$ (1.1), que $Y$ est connexe, et qu'il est un torseur, dans $E_{/U}$, sous le groupe constant $\Aut(Y)_U$.
    \item[(ii)] Tout objet galoisien se trivialise lui-même.
\end{enumerate}

\vskip .3cm
{
Proposition {\bf 2.3.2.} --- \it Soient $A$ et $B$ deux objets connexes de $E$ tels que $A \times B$ soit non-vide. Si $A$ et $B$ se trivialisent l'un l'autre, ils sont isomorphes.
}
\vskip .3cm

{\it Démonstration}. En effet, il existe alors des ensembles non-vides $I$ et $J$ tels que $I_A$ et $J_B$ soient tous deux isomorphes à $A \times B$. Puisque les composantes connexes de $I_A$ (resp. $J_B$) sont toutes isomorphes à $A$ (resp. $B$), $A$ et $B$ sont isomorphes.

\vskip .3cm
{
Corollaire {\bf 2.3.3.} --- \it 
Soient $A$ et $B$ deux objets galoisiens de $E$. Si $\Hom(A, B)$ et $\Hom(B, A)$ sont non-vides, alors $A$ et $B$ sont isomorphes.
}

\vskip .3cm
{
Proposition {\bf 2.3.4.} --- \it Soient $A$ et $B$ deux objets galoisiens de $E$. Si $A$ et $B$ sont isomorphes, tout morphisme de $A$ dans $B$ est un isomorphisme.
}
\vskip .3cm

{\it Démonstration}. Il suffit de prouver que tout morphisme $A \to A$ est un automorphisme. Or (2.1.3.) le $A$-morphisme canonique
$$
\Hom(A, A)_A \to A \times A
$$
est un isomorphisme, et sa restriction à $\Aut(A)_A$ est un isomorphisme puisque $A$ est un pseudo-torseur sous $\Aut(A)_A$. Donc
$$
\Hom(A, A) = \Aut(A).
$$

\vskip .3cm
{\bf 2.3.5.} Soient $G$ un groupe, $B^G$ le topos des $G$--ensembles à droite et $T^G \in \Ob(B^G)$ l'ensemble $G$ muni de l'opération de $G$ par translations à droite.

L'opération de $G$ sur l'ensemble $G$ par translations à gauche fait de $T^G$ un torseur de $B^G$ sous le groupe constant $G_{B^G}$. Pour tout morphisme de topos $E \xrightarrow{u} B^G$, on a donc une structure de $G_E$-torseur à gauche sur $u^{-1}(T^G)$ ; d'où un foncteur
$$
\Homtop(E, B^G)^\circ \to \Tors(E, G_E) \leqno{\text{(*)}}
$$
de la catégorie opposée à $\Homtop(E, B^G)$ dans la catégorie des $G_E$-torseurs à gauche de $E$.
\vskip .3cm
{ 
Lemme. --- \it Le foncteur (*) est une équivalence de catégories.
}
\vskip .3cm
{\it Démonstration abrégée}.
\begin{enumerate}
    \item[a)] Le foncteur (*) est pleinement fidèle parce que $\{T^G\}$ engendre $B^G$.
    \item[b)] Pour tout objet $F$ de $B^G$, l'opération de $G$ sur l'ensemble $F$ fournit une opération du groupe constant $G_E$ sur l'objet constant $F_E$.
\end{enumerate}
Étant donné un $G_E$-torseur à gauche $T$, le foncteur
\begin{align*}
u^{-1}_T: B^G & \to E \\
F & \mapsto T \wedge_{G_E} F_E
\end{align*}
(produit contracté) définit un morphisme de topos $u_T: E \to B^G$ et $u^{-1}_T(T^G)$ est isomorphe à $T$ en tant que $G_E$-torseur.

\vskip .3cm
{\bf 2.3.6.} (``Théorie de Galois'')

Nous supposons le topos $E$ connexe.

Soient $Y$ un objet galoisien de $E$ et $G$ le groupe des automorphismes de $Y$.

On a une structure de $G_E$-torseur à gauche sur $Y$ (2.3.1., remarque (i)), d'où, suivant 2.3.5., un morphisme de topos
$$
u: E \to B^G
$$
tel que le $G_E$-torseur $u^{-1}(T^G)$ soir isomorphe au $G_E$-torseur $Y$.

Le foncteur image directe $u_*$ est isomorphe au foncteur
$$
\phi: E \to B^G
$$
qui associe à l'objet $X$ de $E$ l'ensemble
$$
\phi(X) = \Hom(Y, X)
$$
muni de l'opération à droite de $G = \Aut(Y)$ par composition (en effet, tout faisceau $F$ sur $B^G$ est canoniquement représenté par l'ensemble $F(T^G)$ muni de l'opération à droite de $G$ déduite de l'opération de $G$ sur l'objet $T^G$).

Soit $\LC(E, Y)$ la sous-catégorie pleine de $E$ formée des objets qui sont trivialisés par $Y$ (N. B. ces objets sont localement constants puisque $Y$ recouvre l'objet final). Nous allons prouver que
\begin{enumerate}
    \item[(i)] Le foncteur image inverse
    $$
    u^{-1}: B^G \to E
    $$
    est pleinement fidèle et prend ses valeurs dans $\LC(E, Y)$.
    \item[(ii)] La restriction à $\LC(E, Y)$ du foncteur image directe $u_*\simeq\phi$ est pleinement fidèle.
    D'où les corollaires~:
    \item[(iii)] Les foncteurs $\phi_{/\LC(E, Y)}$ et $u^{-1}$ fournissent des équivalences quasi-inverses
    $$
    \xymatrix{ \LC(E,Y) \ar@<-1ex>[r]  & \ar@<-1ex>[l] B^G}
    $$
    \item[(iv)] La catégorie $\LC(E, Y)$ est un $\cU$-topos et l'inclusion $\LC(E, Y) \to E$ définit un morphisme de topos $E \to \LC(E, Y)$.
\end{enumerate}

{\it Démonstration}.
\begin{enumerate}
    \item[(i)] $T^G$ trivialise tous les objets de $B^G$; donc, étant donné un objet $F$ de $B^G$, $u^{-1}(T^G) \isom Y$ trivialise $u^{-1}(F)$. Prouvons maintenant que le morphisme canonique
    $$
    F \to \phi u^{-1}(F)
    $$
    donné par l'adjonction entre $\phi$ et $u^{-1}$ est un isomorphisme (ce qui entraîne la pleine fidélité de $u^{-1}$).
    \begin{enumerate}
        \item[a)] Le morphisme
        $$
        T^G \to \phi u^{-1}(T^G)
        $$
        est un isomorphisme. En effet, $\phi u^{-1}(T^G)$ est un $G$-torseur puisque $u^{-1}(T^G)$ est isomorphe à $Y$ (2.3.4).
        \item[b)] $\phi$ commute aux sommes directes. (1.2)
        \item[c)] Puisque $T^G$ trivialise $F$, le morphisme
        $$
        T^G \times F \to \phi u^{-1}(T^G \times F) \isom \phi u^{-1}(T^G) \times \phi u^{-1}(F)
        $$
        est un isomorphisme par (a) et (b) ; donc il en est de même du morphisme
        $$
        F \to \phi u^{-1}(F)
        $$
        puisque $T^G$ recouvre l'objet final de $B^G$.
    \end{enumerate}
    \item[(ii)] Soit $L$ un objet de $\LC(E, Y)$. Le $Y$-isomorphisme :
    $$
    \phi(L)_Y = \Hom(Y, L)_Y \isommap Y \times L \leqno{(2.1.3.)}
    $$
    montre qu'on obtient pour tout objet $M$ de $E$ une bijection : 
    $$
    \Hom(Y \times L, M) \isommap \App(\phi(L), \phi(M))
    $$
    en associant au morphisme $m: Y \times L \to M$ l'application
    \begin{align*}
    \phi(L) &\to \phi(M)\\
    f & \mapsto m(f) = m \circ (1_Y, f).
    \end{align*}
    Or on a le diagramme commutatif
    $$
    \xymatrix{
        \Hom(L,M) \ar@{_{(}->}[d] \ar[dr] \\
        \Hom(Y\times L,M) \ar[r]^-\sim & \App(\phi(L), \phi(M))
    }
    $$
    (où la flèche verticale désigne la composition avec la projection $Y \times L \to L$). L'application
    $$
    \Hom(L, M) \to \App(\phi(L), \phi(M))
    $$
    est donc injective ; il reste seulement à prouver que son image est formée des applications qui sont de morphismes de $G$-ensembles.
\end{enumerate}

\vskip .3cm
{\bf 2.4.} Nous supposons à nouveau le topos $E$ localement connexe. Et nous %f
faisons l'hypothèse suivante sur l'univers $\cU$ :

$\cU$ admet un élément de cardinal infini.

Désignons par $\SLC(E)$ la sous-catégorie pleine de $E$ formée des objets qui sont sommes directes d'objets localement constants (i.e. des objets dont les composantes connexes sont des objets localement constants 2.2.5). Les points 2.4.1. à 2.4.10 qui suivent vont prouver le

\vskip .3cm
{
Théorème {\bf 2.4.} --- \it 
\begin{enumerate}
    \item[(i)] La catégorie $\SLC(E)$ est un $\cU$-topos, et l'inclusion $\SLC(E) \to E$ définit un morphisme de topos $E \to \SLC(E)$.
    \item[(ii)] Les objets galoisiens de $\SLC(E)$, qui s'identifient à ceux de $E$, engendrent le topos $\SLC(E)$.
\end{enumerate}
}

\vskip .3cm
{
Proposition {\bf 2.4.1.} --- \it
\begin{enumerate}
    \item[(i)] $\SLC(E)$ est stable dans $E$ par sommes directes et limites projectives finies.
    \item[(ii)] Étant donnés un objet $A$ de $\SLC(E)$ et une relation d'équivalence $R \hookrightarrow A \times A$ de $E$, le quotient $A/R$ est dans $\SLC(E)$ dès que $R$ y est. 
\end{enumerate}
}
\vskip .3cm

Compte tenu de (2.2.4), (2.2.5) et (2.2.6, a), cela découle du lemme suivant :

\vskip .3cm
{
Lemme {\bf 2.4.2.} --- \it Soient $T$ un $\cU$-topos, $K$ une sous-catégorie pleine de $T$, $S$ la sous-catégorie pleine de $T$ formée des objets qui sont sommes directes d'objets de $K$. Supposons que:
\begin{enumerate}
    \item[1)] Tout objet de $K$ est somme directe d'objets connexes $\in \Ob(K)$ ; et tout sous-objet d'un objet $X$ de $K$ qui est somme directe de composantes connexes de $X$ est dans $K$.
    \item[2)] Tout morphisme d'un objet non-vide $\in \Ob(K)$ dans un objet connexe $\in \Ob(K)$ est un épimorphisme (dans $T$).
    \item[3)] $K$ est stable dans $T$ par limites projectives finies.
    \item[4)] Pour tout objet $X$ de $K$ et toute relation d'équivalence $R \hookrightarrow X \times X$ telle que $R \in \Ob(K)$, le quotient $X_{|_R}$ est dans $K$.
\end{enumerate}
Alors :
\begin{enumerate}
    \item[(i)] $S$ est stable dans $E$ par sommes directes et par limites projectives finies.
    \item[(ii)] Pour tout objet $A$ de $S$ et toute relation d'équivalence $R \hookrightarrow A \times A$ telle que $R \in \Ob(S)$, le quotient $A_{|_R}$ est dans $S$.
\end{enumerate}
}
\vskip .3cm

{\it Démonstration}. Le point (i) est immédiat ; prouvons (ii). Nous considérons un diagramme cartésien et cocartésien de $T$
$$
    \xymatrix{
        R \ar[r]^{p_1} \ar[d]_{p_2} & A \ar[d]^\pi\\
        A \ar[r]_\pi & Q
    }
$$
où $A$ et $R$ sont objets de $S$. D'après (1), $A$ est somme directe d'objets connexes $\in \Ob(K)$
$$
\amalg_{i \in I} C_i \isommap A.
$$
Pour tout couple d'indices $(i, j)$, soit
$$
R_{i j} = C_i \times_Q C_j.
$$
Prouvons que $R_{i j}$ est objet de $K$. $R_{i j}$ est produit fibré de $R$ avec $C_i \times C_j$ au-dessus de $A \times A$ ; puisque ces trois objets sont dans $S$, $R_{i j}$ est dans $S$ ; donc $R_{i j} \hookrightarrow C_i \times C_j$ est somme directe de composantes connexes de $C_i \times C_j$, qui est dans $K$ ; donc par (1), $R_{i j}$ est dans $K$.

Pour tout $i \in I$, soit $M_i \hookrightarrow Q$ l'image de $\pi_{|_{C_i}}$. On a $C_i \times_{M_i} C_i = C_i \times_Q C_i=  R_{i i}$, donc le diagramme
$$
    \xymatrix{
        R_{ii} \ar[r]^{p_1} \ar[d]_{p_2} & C_i \ar[d]^\pi\\
        C_i \ar[r]_{\pi} & M_i
    }
$$
est cartésien et cocartésien ; puisque $R_{i i}$ et $C_i$ sont dans $K$, $M_i$ est dans $K$ (par (4)).

Soient $i, j \in I$. Si $R_{i j}$ est vide, alors $M_i \times_Q M_j$ est vide ; si $R_{i j}$ n'est pas vide, les projections
$$
p_1: R_{i j} \to C_i \quad p_2: R_{i j} \to C_j
$$
sont des épimorphismes par (2), donc $M_i$ et $M_j$ définissent le même sous-objet de $Q$. Puisque les $M_i$ sont dans $K$ et recouvrent $Q$, cela prouve que $Q$ est dans $S$, c.q.f.d.

\vskip .3cm
{
Corollaire (de 2.4.2) {\bf 2.4.3.} --- \it 
S'il existe en outre une petite famille $(Y_{\alpha})$ d'objets de $K$ telle que tout objet de $K$ puisse être recouvert (dans $T$) par des $Y_{\alpha}$, alors $S$ est un $\cU$-topos et l'inclusion $S \to T$ définit un morphisme de topos $T \to S$ (par le critère de Giraud).
}

\vskip .3cm
{
Proposition {\bf 2.4.4.} --- \it Soient $L$ et $M$ deux objets de $E$. Tout objet connexe de $E$ qui trivialise $L$ et $M$ trivialise le faisceau $\underline{\Isom}(L, M)$ des germes d'isomorphisme de $L$ dans $M$. (immédiat)
}
\vskip .3cm 

\vskip .3cm
{\bf 2.4.5.}
\begin{enumerate}
    \item[(i)] {Définition. --- \it Soient $G$ un groupe de $E$, $F$ un objet de $E$ et $G \times F \xlongrightarrow{a} F$ une opération de $G$ sur $F$. On dit que $G$ opère \emph{transitivement} sur $F$ si les conditions équivalentes que voici sont remplies~:
    \begin{enumerate}
        \item[1)] Le quotient $F/G$ est un sous-objet de l'objet final (un ouvert).
        \item[2)] Le morphisme $(a, pr_2): G \times F \to F \times F$ est un épimorphisme.
    \end{enumerate}}
    \item[(ii)] {Lemme. --- \it Soit $L$ un objet localement constant de $E$. Si le groupe constant $\Aut(L)_E$ opère transitivement sur $L$, les composantes connexes de $L$ sont des objets \emph{galoisiens} ; et deux composantes connexes de $L$ situées au-dessus de la même composante connexe de $e_E$ sont toujours \emph{isomorphes}.}
\end{enumerate}
\vskip .3cm

{\it Démonstration}. Soit $K$ une composante connexe de $L$.
\begin{enumerate}
    \item[a)] Pour tout $\alpha \in \Aut(L)$, l'image $\alpha (K)$ de $\alpha_{/K}$ est une composante connexe de $L$;
    \item[b)] Soit $H$ le sous-groupe de $\Aut(L)$ formé des automorphismes $\alpha$ tels que $\alpha(K) = K$ comme sous-objets de $L$. D'après (a) on a un diagramme cartésien :
    $$
    \xymatrix{
        H_E \times K \ar@{_{(}->}[d] \ar[r] & K\times K \ar@{_{(}->}[d] \\
        \Aut(L)_E\times L \ar[r] & L\times L
    }
    $$
    et ainsi $H_E$ opère transitivement sur $K$, donc aussi $\Aut(K)_E$. Mais d'après (2.2.6 (b)), $\Aut(K)_E \times K \to K \times K$ est un monomorphisme, donc $K$ est galoisien. Le reste de la proposition est trivial.
\end{enumerate}

\vskip .3cm
{
Proposition {\bf 2.4.6.} --- \it Soit $L$ un objet localement constant de $E$. Au-dessus de chaque composante connexe de $e_E$, il existe un objet galoisien $Y$ qui trivialise $L$ et tel que tout objet connexe de $E$ qui trivialise $L$ trivialise $Y$.
}
\vskip .3cm

{\it Démonstration}. Il suffit de le prouver pour $E$ connexe. Soit $U$ un objet connexe non-vide qui trivialise $L$ et soit $I = \Hom(U, L)$. Représentons le faisceau $\Isom(I_E, L)$ par un objet $T$ de $E$, d'où un $T$-isomorphisme
$$
I_T \isom T \times I_E \to T \times L.
$$
Tout objet de $E$ qui trivialise $L$ trivialise $T$ (2.4.4.) et $T$ est un pseudo-torseur (en fait un torseur 2.2.6. a)) sous le groupe constant des permutations de $I$ ; d'où notre proposition, par 2.4.5 (ii) et 2.2.5 (ii).

\vskip .3cm
{
Corollaire {\bf 2.4.7.} --- \it Pour tout objet $X$ de $\SLC(E)$, il existe des objets galoisiens $(Y_{\alpha})$ de $E$ et une famille épimorphique $(Y_{\alpha} \to X)$ de $E$.
}

\vskip .3cm
{
Proposition {\bf 2.4.8.} --- \it Soient $Y$ un objet connexe de $E$ et $U$ l'image de $Y \to e_E$. Pour que $Y$ soit galoisien, il faut et il suffit qu'il se trivialise lui-même et que $U$ soit une composante connexe de $e_E$.
}
\vskip .3cm

{\it Démonstration}. On sait déjà que ces conditions sont nécessaires (2.3.1, Remarque (i)). Prouvons qu'elles suffisent. Soit $V$ un sous-objet de $e_E$ tel que $U \amalg V \isom e_E$. $Y$ et $V$ trivialisent $Y$, donc $Y$ est localement constant ; puisque en outre $Y$ est connexe, il existe un objet galoisien $Z$ au-dessus de $U$, qui trivialise $Y$ et est trivialisé par $Y$ (2.4.6). Le produit $Y \times Z$ n'est pas vide puisque $U$ ne l'est pas, donc $Y$ et $Z$ sont isomorphes (2.3.2), c.q.f.d.

\vskip .3cm
{
Proposition {\bf 2.4.9.} --- \it La catégorie des objets galoisiens de $E$ est $\cU$-petite à équivalence près.
}
\vskip .3cm

{\it Démonstration}. Compte tenu de l'hypothèse sur l'univers $\cU$, cela résulte des deux points suivants :
\begin{enumerate}
    \item[(i)] Soient $S$ une catégorie $\cU$-petite et $I$ un ensemble $\cU$-petit. L'ensemble des préfaisceaux $F$ sur $S$ tels que $F(X) \subset I$ pour tout $X \in \Ob(S)$ est $\cU$-petit. (Clair)
    \item[(ii)] Soient $G \subset \Ob(E)$ une petite famille génératrice de $E$, formée d'objets connexes non-vides. Soient $c$ le cardinal de l'ensemble des flèches entre objets $\in G$, et $d$ le cardinal dénombrable. Pour tout objet galoisien $Y$ de $E$ et tout $U \in \Ob(G)$, on a :
    $$
    \Card(\Hom(U, Y)) \leq c^d.
    $$
    
    Démonstration. Soit $A$ l'ensemble des flèches $U \to Y$, $U$ parcourant $G$. Étant donné un élément $U \xlongrightarrow{s} Y$ de $A$, l'ensemble des $(V \xlongrightarrow{t} Y) \in A$ tels que $U \times_Y V$ soit non-vide a son cardinal au plus égal à $c^2$ :
    
    Étant données des flèches $W \to U$, $W \to V$ avec $V$, $W \in G$, il existe au plus une flèche $t: V \to Y$ qui rend commutatif le diagramme
    $$
    \xymatrix{
        W\ar[r] \ar[d] & V \ar[d]^t\\
        U \ar[r]_s & Y
    }
    $$
    (en effet, ou bien $\Hom(V, Y) = \emptyset$, ou bien $V$ trivialise $Y$ et alors $\Hom(V, Y) \to \Hom(W, Y)$ est bijective). Puisque l'ensemble des flèches appartenant à $A$ recouvre $Y$, on déduit de (1.3) que :
    $$
    \Card(A) \leq \sum_{n \geq 0} c^{2n} \leq c^d
    $$
    d'où a fortiori notre assertion.
\end{enumerate}

\vskip .3cm
{
Proposition {\bf 2.4.10.} --- \it Soit $u: A \to B$ un morphisme de topos tel que le foncteur $u^{-1}$ soit \emph{pleinement fidèle}.
\begin{enumerate}
    \item[(i)] Pour qu'un objet $Y$ de $B$ soit connexe et non-vide, il faut et il suffit que $u^{-1}(Y)$ le soit.
    \item[(ii)] Supposons $B$ localement connexe. Pour qu'un objet $Y$ de $B$ soit galoisien, il faut et il suffit que $u^{-1}(Y)$ le soit.
\end{enumerate}
}
\vskip .3cm

{\it Démonstration}.
\begin{enumerate}
    \item[(i)] On a pour tout ensemble $I$ des bijections naturelles
    $$
    \Hom(Y, I_B) \isommap \Hom(u^{-1}(Y), u^{-1}(I_B)) \isommap \Hom(u^{-1}(Y), I_A). \quad \text{(cf 1.2)}
    $$
    \item[(ii)]
    \begin{enumerate}
        \item[a)] L foncteur $u_*$ commute aux sommes directes : en effet, pour tout objet connexe non-vide $D$ de $B$, le foncteur $\Hom_A(u^{-1}(D), -)$ y commute (i). 
        \item[b)] Supposons $u^{-1}(Y)$ galoisien. Soit $V$ l'image de $Y \to e_B$. $u^{-1}(V)$ est connexe et non-vide, et il existe un objet $U'$ de $A$ tel que 
        $$
        u^{-1}(V) \amalg U' \isom e_A
        $$
        on a alors (d'après (a))
        $$
        u_* u^{-1}(V) \amalg u_*(U') \isom u_*(e_A) \isom e_B.
        $$
        Or $u_*u^{-1}(V)$ est isomorphe à $V$ puisque $u^{-1}$ est pleinement fidèle. Donc $V$ est une composante connexe de $e_B$ d'après (i).
        
        $Y$ se trivialise lui-même puisque $u^{-1}$ se trivialise lui-même et que $u^{-1}$ est pleinement fidèle ; $Y$ est connexe et non-vide d'après (i). Donc $Y$ est galoisien d'après 2.4.8.
        \item[c)] Réciproquement, supposons $Y$ galoisien. $u^{-1}(Y)$ est alors localement constant, connexe et non-vide (i) ; et puisque le morphisme de groupes $\Aut(Y) \to \Aut(u^{-1}(Y))$ donné par le foncteur $u^{-1}$ est un isomorphisme, $u^{-1}$ est un pseudo-torseur sous le groupe constant $\Aut(u^{-1}(Y))_A$.
    \end{enumerate}
\end{enumerate}

\vskip .3cm
{\bf 2.4.11.} Le théorème 2.4 est maintenant prouvé : 

Le point (i) par 2.4.1., 2.4.7., 2.4.9. et le critère de Giraud ;

D'après 2.4.7. et 2.4.10, (i), $\SLC(E)$ est localement connexe ; d'où le point (ii) par 2.4.10 (ii).

\vskip .3cm
{
Proposition {\bf 2.4.12.} --- \it Soient $F$ un second $\cU$-topos localement connexe et $E \to F$ un morphisme. Le foncteur image inverse transforme les objets de $\SLC(F)$ en objets de $\SLC(E)$, d'où un morphisme de topos $\SLC(E) \to \SLC(F)$.
}


% End

% Begin


%%%%%%%%%%%%%%%%%%%%%%%%%%%%%%%%%%%%%%%%%%%%%%%%%%%%%%%%%%%%%%%
\chapter*{\S \space III. --- TOPOS LOCALEMENT GALOISIENS ET GROUPOÏDE FONDAMENTAL}
\addcontentsline{toc}{section}{III. Topos localement galoisiens et groupoïde fondamental}
\section*{}

On maintient l'hypothèse (2.4) que l'univers $\cU$ admet un élément de cardinal infini.
\vskip .3cm
{
Définition {\bf 3.1.} --- \it Nous dirons qu'un topos $E$ est localement galoisien s'il est engendré par ses objets galoisiens (2.3). Suivant (2.4), cela revient à dire que $E$ est localement connexe et que :
$$
\SLC(E) = E.
$$
}
Dans les numéros (3.2) et (3.3) on montre que les topos localement galoisiens se comportent comme des ``pro-groupoïdes''. On passe ensuite à la définition du (pro-) groupoïde fondamental (3.4).
\vskip .3cm
{\bf 3.2.} (Les groupoïdes comme topos). Soient $\Top$ la 2-catégorie des $\cU$-topos et $\Grpd$ la 2-catégorie des groupoïdes qui sont $\cU$-petits à équivalence près ; soit enfin $\cG$ la sous-catégorie pleine de $\Top$ formée des topos dont tous les objets sont localement constants (ces topos sont donc a fortiori localement galoisiens). Nous allons établir des équivalences
$$
\xymatrix@C=4em{\Grpd \ar@<1ex>[r] \ar@{<-}@<-1ex>[r]& \cG}
$$
\vskip .3cm
{\bf 3.2.1.} Le 2-foncteur \hspace{1.5em} $\xymatrix@R=0pt{\Grpd \ar[r] & \Top\\ C \ar@{|->}[r]& \widehat{C}}$\quad. 

Pour tout objet $C$ de $\Grpd$, la catégorie $\widehat{C}$ des $\cU$-préfaisceaux sur $C$ est un $\cU$-topos. La construction du topos $\widehat{C}$ est 2-fonctorielle en $C$ :
\begin{enumerate}
    \item[(a)] A tout foncteur $C \xlongrightarrow{m} D$ entre objets de $\Grpd$ correspond un morphisme de topos
    $$
    \widehat{m}: \widehat{C} \to \widehat{D}
    $$
    défini par le foncteur image inverse
    $$
    \widehat{m}^{\;-1}: \widehat{D} \to \widehat{C}
    $$
    $$
    G \mapsto G \circ m
    $$
    \item[(b)] A tout morphisme de morphismes de $\Grpd$ :
    $$
    \xymatrix{
        C \ar@<2ex>[r]_{\ }="s"^m \ar@<-2ex>[r]^{\ }="t"_n & D
        \ar"s";"t"^\phi
    }
    $$
    correspond un morphisme de foncteurs $\widehat{D} \times C^\circ \to \Ens$ :
    $$
    G(n(X)) \to G(m(X))
    $$
    d'où un morphisme fonctoriel
    $$
    \widehat{n}^{\;-1} \to \widehat{m}^{\;-1}
    $$
    c'est-à-dire un morphisme de morphismes de topos
    $$
    \widehat{m} \to \widehat{n}
    $$
    \item[(c)] Et la compatibilité de ces données avec les diverses opérations de composition se vérifie immédiatement.
\end{enumerate}

\vskip .3cm
{
Proposition {\bf 3.2.2.} --- \it Pour tout objet $C$ de $\Grpd$, le topos $\widehat{C}$ est objet de $\cG$ : tout préfaisceau représentable sur $C$
\begin{enumerate}
    \item[(i)] est connexe et non-vide dans $\widehat{C}$ (clair) 
    \item[(ii)] trivialise tous les objets de $C$ (2.2.1). 
\end{enumerate}
}
2.2.1.

\vskip .3cm
{
Proposition {\bf 3.2.3.} --- \it Soit $C$ un objet de $\Grpd$. Tout foncteur fibre de $\widehat{C}$ est représentable par un objet de $C$, autrement dit, est isomorphe à un foncteur de la forme $F \to F(X)$, où $X$ est un objet de $C$.
}
\vskip .3cm

{\it Démonstration}. Soit $\phi: \widehat{C} \to \Ens$ un foncteur libre de $\widehat{C}$. Il existe un préfaisceau représentable $X$ sur $C$ tel que $\phi(X) \neq \emptyset$. D'après (3.2.2. (ii)) et (2.1.4), $\phi$ est isomorphe à $F \to F(X)$.

\vskip .3cm
{
Proposition {\bf 3.2.4.} --- \it L'équivalence naturelle $C \to \Point(\widehat{C})$. Soit $C$ un objet de $\Grpd$. A tout objet $X$ de $C$, associons le point $p_X$ de $\widehat{C}$ défini par le foncteur fibre $F \to F(X)$ de $\widehat{C}$ ; cela nous donne un foncteur de $C$ dans la catégorie des points de $\widehat{C}$ :
\begin{align*}
C & \to \Point(\widehat{C})\\
X& \mapsto p_X
\end{align*}
On \c{c}ait que ce foncteur est pleinement fidèle ; donc c'est une équivalence de catégories d'après (3.2.3). 
}

\vskip .3cm
{
Proposition {\bf 3.2.5.} --- \it Soit $E$ un topos objet de $\cG$. La sous-catégorie pleine $C$ de $E$ formée des objets galoisiens qui trivialisent tous les objets de $E$ est un objet de $\Grpd$, et elle engendre $E$ ; d'où une équivalence $\widehat{C} \isom E$.
}
\vskip .3cm

{\it Démonstration}. Pour tout composante connexe $U$ de $e_E$, $E_{/U}$ est objet de $\cG$ (2.2.5). On peut donc supposer $E$ connexe.

D'après (2.4.9), il existe une petite famille $(Y_\alpha)$ d'objets galoisiens de $E$ telle que tout objet galoisien de $E$ soit isomorphe à un $Y_\alpha$. Soit $S$ la somme directe des $Y_\alpha$. Tous les objets de $E$ sont localement constants, donc il existe (2.4.6) un objet galoisien $Z$ de $E$ qui trivialise $S$, et de ce fait tous les $Y_\alpha$ (2.2.5). Puisque tout objet de $E$ est trivialisé par un $Y_\alpha$ (2.4.6), $Z$ trivialise tous les objets de $E$ ; et $Z$ engendre $E$ d'après (2.2.6 (a)). A fortiori, $C$ engendre $E$. Enfin, tous les objets de $C$ trivialisent $Z$ et $Z$ trivialise tous les objets de $C$, donc tous les objets de $C$ sont isomorphes (2.3.2) et $C$ est un groupoïde (2.3.4).

\vskip .3cm
{
Corollaire {\bf 3.2.6.} --- \it La catégorie $\Point(E)$ des points de $E$ est un objet de $\Grpd$ (3.2.4.).
}

\vskip .3cm
{
Corollaire {\bf 3.2.7.} --- \it L'équivalence naturelle $E\to (\Point(E))\widehat{\phantom{C}}$
\begin{enumerate}
    \item[a)] Soit $E$ un objet de $\cG$. On définit un \emph{foncteur}
    $$
    E \to (\Point(E))\widehat{\phantom{C}} \leqno{\text{(*)}}
    $$
    en associant à l'objet $X$ de $E$ le préfaisceau
    $$
    p \mapsto p^{-1}(X)
    $$
    sur $\Point(E)$ (l'action sur les morphismes est évidente).
    \item[b)] Le foncteur (*) est une équivalence de catégories (donc il définit une équivalence de topos $(\Point(E))\widehat{\phantom{C}} \to E$).
\end{enumerate}
}
\vskip .3cm

D'après (3.2.5), il suffit de le prouver pour $E = \widehat{C}$, où $C$ est un objet de $\Grpd$. L'équivalence naturelle $C \to \Point(\widehat{C})$ fournit une équivalence de catégories
$$
(\Point(\widehat{C}))\widehat{\phantom{C}} \to \widehat{C}
$$
et le foncteur composé
$$
\widehat{C} \to (\Point(\widehat{C}))\widehat{\phantom{C}} \to \widehat{C}
$$
n'est autre que le foncteur identique de $\widehat{C}$.

\vskip .3cm
{\bf 3.2.8.} {\it Conclusion}. Les foncteurs
\begin{align*}
C \mapsto \widehat{C}\\
\Point(E) \mapsfrom E
\end{align*}
définissent des équivalences quasi-inverses entre $\Grpd$ et $\cG$.

\vskip .3cm
{\bf 3.3.} (Les topos localement galoisiens comme limites projectives filtrantes de groupoïdes)

\vskip .3cm
{\bf 3.3.0.} Dans tout ce numéro 3.3, on entend par ``ordonnés filtrants'' les ensembles ordonnés filtrants à gauche $\cU$-petits, que l'on regarde aussi bien comme des catégories.

\vskip .3cm
{\bf 3.3.1.} Notre propos est d'établir pour tout $\cU$-topos $E$ l'équivalence des propriétés suivantes :
\begin{enumerate}
    \item[(i)] $E$ est localement galoisien.
    \item[(ii)] Il existe un ordonné filtrant $I$ (3.3.0) et une catégorie fibrée en $\cU$-topos $F \to I$ (A. 1) qui remplit les conditions suivantes :
    \begin{enumerate}
        \item[1)] Les fibres de $F$ sont des objets de $\cG$ (3.2).
        \item[2)] Les foncteurs changement de base de $F$ sont pleinement fidèles.
        \item[3)] $E$ est une 2-famille projective de $F$ dans la 2-catégorie des $\cU$-topos (A. 4)
    \end{enumerate}
\end{enumerate}

\vskip .3cm
{
Proposition {\bf 3.3.2.} --- \it Soit $E$ un $\cU$-topos localement connexe. Pour tout crible $R$ couvrant $e_E$, soit $\LC(E, R)$ la sous-catégorie pleine de $E$ formée des objets qui sont trivialisés par les objets connexes appartenant à $R$.
\begin{enumerate}
    \item[(i)] La catégorie $\LC(E, R)$ est un $\cU$-topos objet de $\cG$.
    \item[(ii)] Les inclusions $\LC(E, R) \to \SLC(E)$ et $\LC(E, R) \to E$ définissent des morphismes de topos $\SLC(E) \to \LC(E, R)$ et $E \to \LC(E, R)$.
\end{enumerate}
}
\vskip .3cm

{\it Démonstration}. 
\begin{enumerate}
    \item[1)] $\LC(E, R)$ est stable dans $E$ par limites inductives et projectives (2.2.4) (elle l'est donc aussi dans $\SLC(E)$ par (2.4)).
    \item[2)] Soit $K$ la catégorie des objets de $\LC(E, R)$ qui sont galoisiens dans $E$. D'après (2.4.6) et le point (1), $K$ engendre $\LC(E, R)$. Or d'après (2.4.9), $K$ est équivalente à une petite catégorie, donc $\LC(E, R)$ est un $\cU$-topos et on a le point (ii). En particulier, les objets de $K$ qui trivialisent \emph{dans $E$} un objet donné de $\LC(E, R)$ le trivialisent aussi dans $\LC(E, R)$. Donc tout objet de $\LC(E, R)$ est localement constant (2.4.6).
\end{enumerate}

\vskip .3cm
{
Proposition {\bf 3.3.3.} --- \it Soient $T$ un $\cU$-topos, $I$ un ordonné filtrant (3.3.0) et $(T_i)_{i \in I}$ une famille décroissante de sous-catégories essentiellement pleines de $T$. Soit $F$ la sous-catégorie pleine de $I \times T$ formée des couples $(i, X)$ tels que $X \in \Ob(T_i)$. Supposons qu'on ait les propriétés suivantes :
\begin{enumerate}
    \item[1)] Pour tout $i \in I$, $T_i$ est un $\cU$-topos objet de $\cG$, et l'inclusion $T_i \to T$ définit un morphisme de topos $T \to T_i$. 
    \item[2)] La réunion des $\Ob(T_i)$ engendre $T$. Alors $F$ est une catégorie fibrée en topos au-dessus de $I$ (cf. A 1), et l'inclusion $F \to I \times T$ définit un morphisme de catégories fibrées en topos $I \times T \to F$ qui fait de $T$ la 2-limite projective des $T_i$ dans la catégorie des $\cU$-topos.
\end{enumerate}
}
\vskip .3cm

{\it Démonstration}. Il est immédiat que $F$ est une catégorie fibrée en topos, et que l'inclusion $F \to I \times T$ définit un morphisme de catégories fibrées en topos $I \times T \to F$.

Soit $G$ la sous-catégorie pleine de $T$, réunion des $T_i$. On a tout de suite :
\begin{enumerate}
    \item[(i)] Les objets de $G$ sont localement constants dans $T$.
    \item[(ii)] $G$ est stable par limites projectives finies dans $T$.
\end{enumerate}

Soit maintenant $S$ un $\cU$-topos. Il faut prouver que le foncteur
$$
\phi: \Homtop(S, T) \to \Cartop_I(I \times S, F)
$$
qui associe au morphisme $f: S \to T$ le morphisme $I \times S \to F$ défini par le foncteur
\begin{align*}
\phi(f)^{-1}: F &\to I \times S\\
(i, X) &\mapsto (i, f^{-1}(X))
\end{align*}
est une équivalence de catégories.
\begin{enumerate}
    \item[(i)] $\phi$ est pleinement fidèle : 
    Soient $f, g: S \to T$ deux morphismes ; étant donné un morphisme de foncteurs cartésiens
    $$
    \lambda: \phi(g)^{-1} \to \phi(f)^{-1}
    $$
    il existe un morphisme de foncteurs et un seul
    $$
    \mu: g^{-1}_{/G} \to f^{-1}_{/G}
    $$
    tel que pour tout objet $(i, X)$ de $F$ on ait~: 
    $$
    \lambda_{(i, X)} = (1_i, \mu_X)
    $$
    Mais puisque $G$ engendre $T$, le foncteur
    \begin{align*}
    \Homtop(S, T) & \to (\Fonct(G, S))^\circ \\
    f & \mapsto f^{-1}_{/G}
    \end{align*}
    est pleinement fidèle, d'où notre assertion.
    \item[(ii)] $\phi$ est essentiellement surjectif : 
    
    Soit un morphisme de catégories fibrées en topos
    $$
    I \times S \to F
    $$
    défini par un $I$-foncteur cartésien
    \begin{align*}
    F &\to I \times S\\
    (i, X) & \mapsto (i, a(i, X))
    \end{align*}
    \begin{enumerate}
        \item[a)] Il existe un foncteur $b: G \to S$ tel que le foncteur $a: F \to S$ soit isomorphe au foncteur $(i, X) \mapsto b(X)$ :
        
        Puisque les changements de base de $F$ sont donnés par les inclusions entre les $G_i$, on peut définir pour tout objet $X$ de $G$ :
        $$
        b(X) = \varprojlim_{X \in \Ob(T_i)} a(i, X)
        $$
    Toutes les projections $b(X) \to a(i, X)$ sont des isomorphismes. Étant donné un morphisme $X \to Y$ de $G$, il existe un morphisme et un seul $b(X) \to b(Y)$ tel que, dès que $X$ et $Y$ sont dans $\Ob(T_i)$, on ait le diagramme commutatif :
    $$
        \xymatrix{
            b(X) \ar[r]^\sim \ar[d] & a(i,X) \ar[d]\\
            b(Y) \ar[r]^\sim & a(i,Y)
        }
    $$
     d'où un foncteur $b: G\to S$ qui remplit clairement les conditions requises.
     \item[b)] Le foncteur $b$ est exact à gauche : en effet, les restrictions de $a$ aux fibres de $F$ sont des foncteurs exacts à gauche, ainsi que les inclusions $T_i \to G$.
     \item[c)] Soit $(X_\alpha \to X)$ une famille de morphismes de $G$ qui est épimorphique dans $T$ ; la famille de morphismes $(b(X_\alpha) \to b(X))$ de $S$ est épimorphique~:
     
     Soit $i$ un indice tel que $X \in \Ob(T_i)$. Décomposons $X$ en ses composantes connexes dans $T_i$ :
     $$
     X = \amalg_{\gamma \in C} Y_\gamma .
     $$


     Les $Y_\gamma$ sont encore connexes dans $T$ (2.4.10). Pour tout $\gamma \in C$, il existe un objet non-vide $Z_\gamma$ de $G$ et un morphisme $Z_\gamma \to Y_\gamma$ tel que $Z_\gamma \to Y_\gamma \to X$ se factorise par un $X_\alpha \to X$. Le morphisme $Z_\gamma \to Y_\gamma$ est un épimorphisme (2.2.6) donc aussi $b(Z_\gamma) \to b(Y_\gamma)$ (puisque $Z_\gamma$ et $Y_\gamma$ sont tous deux dans un même $T_j$). Mais $b(X)$ est somme directe des $b(Y_\gamma)$ ; donc la famille $(b(Z_\gamma) \to b(X))$ est épimorphique, d'où notre assertion.
     \item[d)] Il existe donc un morphisme de topos $f: S \to T$ tel que la restriction $f^{-1}_{|G}$ soit isomorphe à $b$ ; le foncteur
     \begin{align*}
     F & \to I \times S \\
     (i, X) & \mapsto (i, f^{-1}(X))
     \end{align*}
     est alors isomorphe au foncteur
     $$
     (i, X) \mapsto (i, a(i, X)),
     $$
     c.q.f.d.
    \end{enumerate}
\end{enumerate}

\vskip .3cm
{
Corollaire {\bf 3.3.4.} --- \it Soient $E$ un $\cU$-topos localement connexe, $I$ un ordonné filtrant (3.3.0) et $(R_i)_{i \in I}$ une famille décroissante de cribles couvrant $e_E$ (i.e. $R_i$ est plus fin que $R_j$ pour $i \leq j$). Supposons que pour tout objet localement constant $L$ de $E$, il existe un $i \in I$ tel que $L$ soit dans $\LC(E, R_i)$ (3.3.2) ; alors le topos $T = \SLC(E)$ est limite projective des $T_i = \LC(E, R_i)$ au sens de (3.3.3).
}
\vskip .3cm

{\it Remarque}. Puisqu'il existe une telle famille de cribles, l'implication $(i) \Rightarrow (ii)$ de (3.3.1) est prouvée.

\vskip .3cm
{
Proposition {\bf 3.3.5.} --- \it Soient $I$ un ordonné filtrant, et $\Pi: F \to I$ une catégorie fibrée au-dessus de $I$. Si les foncteurs changement de base de $F$ sont pleinement fidèles, il existe une catégorie $K$ et une famille décroissante $(K_i)_{i \in I}$ de sous-catégories essentiellement pleines de $K$ telles que :
\begin{enumerate}
    \item[a)] $K$ soit réunion des $K_i$ ;
    \item[b)] $F$ soit $I$-équivalente à la sous-catégorie pleine $F'$ de $I \times K$ formée des couples $(i, X)$ tels que $X \in \Ob(K_i)$ ($F'$ est alors une sous-catégorie fibrée de $I \times K$).
\end{enumerate}
}
\vskip .3cm

{\it Démonstration}. J'abrège un peu la démonstration, qui est d'une lourde trivialité. 
\begin{enumerate}
    \item[1)] Soit $F_{\cart}$ la catégorie des objets de $F$ et flèches cartésiennes. Pour tout objet $X$ de $F$, soit $I_X$ la catégorie $F_{\cart}/X$. On a un foncteur
    $$
    I_X \to F
    $$
    (oubli de la flèche structurale), d'où un autre foncteur
    $$
    I_X \to I
    $$
    composé de $I_X \to F$ et de $F \to I$.
    \item[2)] Pour tout couple $(X, Y)$ d'objets de $F$, soit
    $$
    I_{XY} = I_{X} \times_I I_Y
    $$
    Cette catégorie est non-vide puisque $I$ est filtrante ; et on a des foncteurs 
    $$
    q^X_{XY} = I_{XY} \xlongrightarrow{pr_1} I_X \to F, \quad q^Y_{XY} = I_{XY} \xlongrightarrow{pr_2} I_Y \to F.
    $$
    \item[3)] \emph{Remarque}. Étant donnés un objet
    $$
    C = (X' \to X, Y' \to Y)
    $$
    de $I_{XY}$ et une flèche $f: X' \to Y'$ au-dessus de $\Pi(X') = \Pi(Y')$, il existe un morphisme de foncteurs cartésiens et un seul
    $$
    \lambda: q^X_{XY} \to q^Y_{XY}
    $$
    tel que $\lambda_C = f$.
    \item[4)] Définition de la catégorie $K$:
    \begin{enumerate}
        \item[a)] Les objets de $K$ sont les objets de $F$ ;
        \item[b)] Étant donnés des objets $X$, $Y$ de $K$, les morphismes $X \to Y$ sont les morphismes de foncteurs cartésiens :
        $$
        q^X_{XY} \to q^Y_{XY}
        $$
        \item[c)] La composition des morphismes se définit de fa\c{c}on évidente à partir de la remarque 3.
    \end{enumerate}
    \item[5)] Un foncteur $\epsilon: F \to K$. 
    
    On prend $\epsilon(X) = X$ pour tout objet $X$ de $F$. Définissons maintenant l'action de $F$ sur les morphismes :
    
    Soit $f: X \to Y$ un morphisme de $F$, de projection $u: i \to j$. Prenons une flèche $u$-cartésiennne $Y' \to Y$ de $F$, d'où un objet
    $$
    C = (X \xlongrightarrow{1_X} X, Y' \to Y)
    $$
    de $I_{XY}$. Le morphisme $\epsilon(f)$ est le seul qui rende commutatif le diagramme
    $$
        \xymatrix{
            X \ar[d]_{\epsilon(f)_C} \ar[dr]^f \\
            Y' \ar[r] & Y
        }
    $$
    (morphisme qui ne dépend pas du choix de $Y'\to Y$).
    \item[6)] Soit $f: X \to Y$ une flèche de $F$. Pour que $\epsilon(f)$ soit un isomorphisme, il faut et il suffit que $f$ soit cartésienne.
    \item[7)] Pour tout $i \in I$, la restriction de $\epsilon$ à la fibre $F_i$ est pleinement fidèle
    \item[8)] Pour tout $i \in I$, soit $K_i$ l'image essentielle de $e_{| F_i}: F_i \to K$. La famille $(K_i)$ est décroissante d'après (6), et $K$ est réunion des $K_i$.
    \item[9)] Le foncteur
    $$
    (\Pi, \epsilon): F \to I \times K
    $$
    est $I$-cartésien et pleinement fidèle ; son image essentielle est la sous-catégorie pleine $F'$ de $I \times K$ formée des $(i, X)$ tels que $X \in \Ob(K_i)$.
\end{enumerate}

\vskip .3cm
{
Proposition {\bf 3.3.6.} --- \it Soient $C$ une catégorie, $I$ un ordonné filtrant (3.3.0) et $(G_i)_{i \in I}$ une famille décroissante de sous-catégories essentiellement pleines de $C$. On suppose que~:
\begin{enumerate}
    \item[(i)] Pour tout $i \in I$, $G_i$ est un $\cU$-topos objet de $\cG$ (cf. 3.2) ;
    \item[(ii)] Pour $i \leq j$, l'inclusion $G_ j \to G_i$ définit un morphisme de topos $G_i \to G_j$ ;
    \item[(iii)] $C$ est réunion des sous-catégories $G_i$. Dans ce cas, si l'on munit $C$ de sa topologie canonique, la catégorie $\widetilde{C}$ des $\cU$-faisceaux sur $C$ est un $\cU$-topos localement galoisien, et les foncteurs pleinement fidèles 
    $$
    G_i \to C \to \widetilde{C} \leqno{\text{(*)}}
    $$
    définissent des morphismes de topos $\widetilde{C} \to G_i$.
\end{enumerate}
}
\vskip .3cm

{
Corollaire. --- \it Sous les hypothèses précédentes, soient $(T_i)$ les images essentielles des foncteurs pleinement fidèles (*) ; $\widetilde{C}$ est limite projective des topos $T_i$ au sens de (3.3.3).
}
\vskip .3cm

{\it Preuve de la proposition}. Les points suivants se prouvent tous immédiatement :
\begin{enumerate}
    \item[1)] Les limites projectives finies de $C$ sont représentables, et les inclusions $G_i \to C$ sont exactes à gauche.
    \item[2)] Toute famille de morphismes d'un $G_i$ qui est épimorphique dans $G_i$ est épimorphique effective universelle dans $C$.
    \item[3)] $C$ admet une petite famille topologiquement génératrice pour la topologie canonique (cela découle de (2) puisque $I$ est petit (3.3.0)).
    \item[4)] Les foncteurs pleinement fidèles
    $$
    G_i \to C \to \widetilde{C}
    $$
    définissent des morphismes de topos $\widetilde{C} \to G_i$ (par (1) et (2)).
    \item[5)] $\widetilde{C}$ est localement galoisien : les objets galoisiens des $G_i$ sont galoisiens dans $\widetilde{C}$ (par (4) et (2.4.10)) et ils engendrent le site $C$ (par (2)), donc le topos $\widetilde{C}$.
\end{enumerate}

\vskip .3cm
{\bf 3.3.7.} Les numéros (3.3.5) et (3.3.6) prouvent l'implication (ii) $\Rightarrow$ (i) de 3.3.1. 

\vskip .3cm
{\bf 3.4.} Le groupoïde fondamental

\vskip .3cm
{
Définition {\bf 3.4.1.} --- \it Nous appellerons (par abus de langage) groupoïde fondamental d'un $\cU$-topos $E$ la donnée d'un $\cU$-topos \emph{localement galoisien} $S$ et d'un morphisme de topos :
$$
p: E \to S
$$
tel que pour tout $\cU$-topos localement galoisien $T$ le foncteur
$$
\Homtop(S, T) \to \Homtop(E, T)
$$
donné par la composition avec $p$ soit une équivalence de catégories.
}
\vskip .3cm

{\it Remarque}. Si $E$ est localement connexe, le morphisme de topos $E \to \SLC(E)$ défini par l'inclusion $\SLC(E) \to E$ (2.4) fait de $\SLC(E)$ un groupoïde fondamental de $E$.

\vskip .3cm
{
(Condition suffisante pour qu'un topos connexe admette un groupoïde fondamental)
\vskip .3cm
Proposition {\bf 3.4.2.} --- \it
Soit $E$ un $\cU$-topos connexe. Supposons qu'il existe une petite famille $(Y_{\alpha})$ d'objets galoisiens de $E$, telle que tout objet de $E$ qui admet une structure de torseur sous un groupe constant puisse être trivialisé par un $Y_{\alpha}$. Soient $K$ la sous-catégorie pleine de $E$ formée des objets qui peuvent être trivialisés par un objet galoisien, et $S$ la sous-catégorie pleine de $E$ formée des sommes directes d'objets de $K$. $S$ est un $\cU$-topos localement galoisien, et l'inclusion $S \to E$ fait de $S$ un groupoïde fondamental de $E$.
}
\vskip .3cm

{\it Démonstration}.
\begin{enumerate}
    \item[1)] Pour toute famille finie $(L_i)_{1 \leq i \leq n}$ d'objets de $K$, il existe un indice $\alpha$ tel que $Y_\alpha$ trivialise chacun des $L_i$ : en effet, prenons pour chaque $i$ un objet galoisien qui trivialise $L_i$. Le produit
    $$
    T = Y_1 \times ... \times Y_n
    $$
    trivialise chacun des $L_i$ et admet une structure de torseur sous le groupe constant
    $$
    (\Aut(Y_1) \times ... \times \Aut(Y_n))_E
    $$
    il existe ainsi un $Y_\alpha$ qui trivialise $T$, et donc chacun des $L_i$.
    \item[2)] $S$ est un $\cU$-topos localement galoisien, et l'inclusion $S \to E$ définit un morphisme de topos $E \to S$. D'après (1) et la théorie de Galois (les points (iii) et (iv) de 2.3.6), la sous-catégorie $K$ de $E$ remplit les conditions (1) à (4) du lemme (2.4.2.) - compte tenu de fait que les topos de la forme $B^G$ sont objets de $\cG$ -. Or tout objet de $K$ peut être recouvrement par des exemplaires d'un $Y_{\alpha}$, donc $S$ est un $\cU$-topos et l'inclusion $S \to E$ définit un morphisme de topos : $p: E \to S$ (2.4.3.). Mais alors les $Y_{\alpha}$ sont galoisiens dans $S$ (par (2.4.10)), donc $S$ est localement galoisien.
    \item[3)] $p: E \to S$ est un groupoïde fondamental de $E$. Soit $T$ un $\cU$-topos localement galoisien. Le foncteur
    $$
    \Homtop(S, T) \to \Homtop(E, T)
    $$
    est évidemment pleinement fidèle. Prouvons qu'il est essentiellement surjectif. Soit $u: E \to T$ un morphisme. Si $Z$ est un objet galoisien de $T$, $u^{-1}(Z)$ admet une structure de torseur sous $(\Aut(Z))_E$ (2.1.5), donc $u^{-1}(Z)$ est dans $S$. Par conséquent, le foncteur $u^{-1}$ se factorise par $S$, c.q.f.d.
\end{enumerate}

\vskip .3cm
{\bf 3.4.3.} {\it Question}. Soient $E$ un $\cU$-topos et $p: E \to S$ un groupoïde fondamental de $E$. Le foncteur $p^{-1}: S \to E$ est-il toujours pleinement fidèle ? Si la réponse à cette question était affirmative, on obtiendrait aisément une condition nécessaire et suffisante pour l'existence d'un groupoïde fondamental.



% End

% Begin


%%%%%%%%%%%%%%%%%%%%%%%%%%%%%%%%%%%%%%%%%%%%%%%%%%%%%%%%%%%%%%%
\chapter*{\S \space IV. --- LIMITES INDUCTIVES DE TOPOS ET THÉORÈME DE VAN KAMPEN}
\addcontentsline{toc}{chapter}{IV. Limites inductives de topos et théorème de Van Kampen}
\section*{}

Soient $I$ une petite catégorie, $F \xlongrightarrow{\Pi} I$ une catégorie fibrée en $\cU$-topos $(A, 1)$, et
$$
L = \Cart_I(I, F)
$$
la catégorie des sections cartésiennes de $F$ au-dessus de $I$.

\vskip .3cm
{\bf 4.1. $L$ comme limite inductive de $F$}

\vskip .3cm
{
Proposition {\bf 4.1.1.} --- \it $L$ est un $\cU$-topos, et le foncteur d'évaluation :
\begin{align*}
I \times L &\to F\\
(i, S) & \mapsto S(i)
\end{align*}
définit un morphisme de catégories fibrées en topos $F \to I \times L$ qui fait de $L$ la 2-limite inductive de $F$ dans la catégorie des $\cU$-topos. (A, 4)
}
\vskip .3cm

{\it Démonstration}. 
\begin{enumerate}
    \item[1)] Les limites projectives finies (resp. les limites inductives) sont représentables dans $L$ ; et pour tout $i \in \Ob(I)$, le foncteur
    \begin{align*}
    L &\to F_i \\
    S &\mapsto S(i)
    \end{align*}
    y commute : c'est immédiat puisque les foncteurs changement de base de $F$ sont exacts à gauche et commutent aux limites inductives.
    \item[2)] $L$ est un $\cU$-topos : compte tenu du critère de Giraud et de (1), il suffit de montrer que $L$ admet une petite famille génératrice. Or (comme par hasard) $F$ remplit les conditions du corollaire I.9.25 de \cite{3}, qui garantissent l'existence d'une telle famille.
    \item[3)] Le foncteur $I \times L \to F$ définit un morphisme de catégories fibrées en topos $F \to I \times L$~: ce foncteur est évidemment cartésien ; d'où notre assertion par (1).
    \item[4)] Le morphisme de catégories fibrées en topos $F \to I \times L$ fait de $L$ la 2-limite inductive de $F$ dans la catégorie des $\cU$-topos : soit $E$ un $\cU$-topos. Le foncteur
    $$
    \Homtop(L, E) \to \Cartop_I(F, I \times E)
    $$
    qui associe au morphisme $u: L \to E$ le morphisme de catégories fibrées en topos
    $$
    F \to I \times E
    $$
    défini par le foncteur
    \begin{align*}
    I \times E &\to F\\
    (i, X) &\mapsto u^{-1}(X)(i)
    \end{align*}
    est une équivalence de catégories, comme on le vérifie tout de suite en tenant compte de (1).
\end{enumerate}
 
\vskip .3cm
{
Proposition {\bf 4.1.2.} --- \it Si pour tout $i \in \Ob(I)$ la fibre $F_i$ est localement connexe, alors $L$ est localement connexe.
}
\vskip .3cm

{\it Démonstration}. Pour tout $i \in \Ob(I)$, soit
$$
c_i: F_i \to \Ens
$$
le foncteur ``composante connexes'' de $F_i$. (1.5)

Soit $S$ une section cartésienne de $F$. Pour toute flèche $u: i \to j$ de $I$, on a une application
$$
c_i(S(i)) \to c_j(S(j))
$$
qui associe à la composante connexe $C$ de $S(i)$ l'unique composante connexe $D$ de $S(j)$ telle que la restriction à $C$ du $u$-morphisme $S(u): S(i) \to S(j)$ se factorise par $D$.

Étant donnés une seconde section cartésienne $T$ de $F$ et un morphisme $S \to T$ on a pour tout flèche $u: i \to j$ de $I$ le diagramme commutatif :
$$
\xymatrix{
    c_i(S(i)) \ar[r] \ar[d] & c_j(S(j)) \ar[d]\\
    c_i(T(i)) \ar[r] & c_j(T(j))
}
$$

On peut donc définir un foncteur
\begin{align*}
c: L &\to \Ens\\
S &\mapsto c(S) = \varinjlim_I c_i(S(i))
\end{align*}
Prouvons que $c$ est un foncteur composantes connexes de $L$ (cf 1.5.) :

Soit $A$ un ensemble. Considérons l'objet constant $A_L$ de $L$. Pour tout $i \in \Ob(I)$, $A_L(i)$ s'identifie à $A_{F_i}$ ; pour tout flèche $u: i \to j$ de $I$ et tout $a \in A$, le diagramme
$$
\xymatrix{
    e_{F_i} \ar[r] \ar[d]_a & e_{F_j} \ar[d]^a \\
    A_{F_i} \ar[r]_{A_L(u)} & A_{F_j}
}
$$
est commutatif.

Soit maintenant $S$ un objet de $L$. Pour tout $i \in \Ob(I)$, soit $f_i: S(i) \to A_{F_i}$ un $i$-morphisme, correspondant à une application $b_i: c_i(S(i)) \to A$. Étant donnée une flèche $u: i \to j$ de $I$, la commutativité du diagramme
$$
\xymatrix{
    S(i) \ar[r] \ar[d] & S(j) \ar[d] \\
    A_{F_i} \ar[r] & A_{F_j}
}
$$
équivaut à celle du diagramme
$$
\xymatrix@C=1em{
c_i(S(i)) \ar[rr] \ar[dr] && \ar[dl] c_j(S(j)) \\
& A
}
$$
Donc le foncteur $c$ est adjoint à gauche du foncteur ``objet constant'' $\Ens \to L$, c.q.f.d.

\vskip .3cm
{\bf 4.2. Le groupoïde fondamental de $L$}

On définit une \emph{topologie} sur la catégorie $F$, qui fournit un topos $\widetilde{F}$, et un morphisme de topos $\widetilde{F} \to L$ (4.4). Quand les fibres de $F$ sont localement connexes, $\widetilde{F}$ est lui-même localement connexe, et les morphismes de topos $\SLC(\widetilde{F}) \to \SLC(E)$ déduit du morphisme $\widetilde{F} \to L$ est une équivalence (th. 4.5). On montre alors (4.6) comment construire une famille convenable de groupoïdes ``approchés'' $\LC(\widetilde{F}, R)$ (au sens de 3.3.2 et 3.3.4) à l'aide de groupoïdes ``approchés'' des fibres de $F$.

\vskip .3cm
{
Proposition {\bf 4.3.} --- \it Un foncteur pleinement fidèle
$$
p: L = \Cart_I(I, F) \to \widehat{F}
$$
(cette construction utilise seulement le fait que $\Pi: F \to I$ est un foncteur fibrant).
\begin{enumerate}
    \item[(1)] Soit $S: I \to F$ une section cartésienne. On définit le préfaisceau $p(S)$ comme voici :
    Pour tout objet $X$ de $F$, on prend :
    $$
    p(S)(X) = \Hom_i(X, S(i)), \quad \text{où} \quad i = \Pi(X)
    $$
    Étant donné un morphisme $f: X \to Y$ de $F$, l'application $p(S)(f): p(S)(Y) \to p(S)(Y)$ est la composée~:
    $$
    \Hom_j(Y, S(j))\to \Hom_u(X, S(j)) \isommap \Hom_i(X, S(i))
    $$ 
    où $u: i \to j$ désigne la projection de $f$. Autrement dit, l'application $p(S)(f)$ envoie l'élément $h$ de $p(S)(Y)$ sur l'élément $g$ de $p(S)(X)$ qui rend commutatif le diagramme
    $$
    \xymatrix{
        X \ar[r]^f \ar[d]_g & Y \ar[d]^h\\
        S(i) \ar[r]_{S(u)} & S(j)
    }
    $$
    Il est donc clair qu'on a bien défini un préfaisceau sur $F$.
    \item[(2)] Soient maintenant $S, T: I \to F$ deux sections cartésiennes et $m: S \to T$ un morphisme. On définit le morphisme
    $$
    p(m): p(S) \to p(T)
    $$
    comme voici : soit $X$ un objet de $F$, de projection $i$. L'application
    $$
    p(m)_X: p(S)(X) \to p(T)(X)
    $$
    est la composition avec le $i$-morphisme :
    $$
    m_i: S(i) \to T(i)
    $$
    fourni par $m$.
    \item[(3)] Le foncteur $p$ est pleinement fidèle.
\end{enumerate}
}
\vskip .3cm

{\it Démonstration}. Soient $S$, $T$ deux sections cartésiennes de $F$ et $\mu: p(S) \to p(T)$ un morphisme. Pour tout $i \in \Ob(I)$, soit
$$
\mu_i: p(S)_{/F_i} \to p(T)_{/F_i}
$$
le morphisme induit par $\mu$ entre les restrictions à la fibre $F_i$. Comme ces restrictions sont représentées par $S(i)$ et $T(i)$ respectivement, $\mu_i$ provient d'un $i$-morphisme
$$
m_i: S(i) \to T(i)
$$
à savoir $m_i = \mu_{S(i)}(1_{S(i)})$

Prouvons que les $m_i$ définissent un morphisme $m: S\to T$ : 

Soit $i \xlongrightarrow{u} j$ une flèche de $I$. On a le diagramme commutatif
$$
    \xymatrix{
        p(S)(S(i)) \ar[r]^\mu & p(T)(S(i))\\
        p(S)(S(j)) \ar[r]_\mu \ar[u] & p(T)(S(j)) \ar[u]
    }
$$
La flèche de gauche transforme $1_{S(j)}$ en $1_{S(i)}$ ; donc la flèche de droite transforme $m_j$ en $m_i$, autrement dit le diagramme
$$
    \xymatrix{
        S(i) \ar[r] \ar[d]_{m_i} & S(j) \ar[d]^{m_j} \\
        T(i) \ar[r] & T(j)
    }
$$
est commutatif, cqfd.

Soit $n: S \to T$ un morphisme. Pour que $p(n) = \mu$, il faut et il suffit que $n = m$, comme on le vérifie aisément. Donc $p$ est pleinement fidèle.

\vskip .3cm
{\bf 4.4. Le topos $\widetilde{F}$}

\vskip .3cm
{ 
Définition (Topologie $T$ sur $F$) {\bf 4.4.1.} --- \it 
Soit $X$ un objet de $F$, de projection $i \in \Ob(I)$. Un crible $R$ de $F_{/X}$ est couvrant pour $T$ s'il contient une famille épimorphique de la fibre $F_i$.

Les axiomes d'une topologie se vérifient à l'aide des remarques suivantes :
\begin{enumerate}
    \item[(i)] Soit $u: i \to j$ une flèche de $I$, et considérons un diagramme commutatif de $F$ :
    $$
    \raisebox{2em}{\xymatrix{
        P \ar[r] \ar[d] & Z' \ar[d] \\
        X \ar[r] & Z
    }}
    \quad \text{(*)}
    $$
    où $P \to X$ est un $i$-morphisme, $Z' \to Z$ un $j$-morphisme, et les flèches horizontales des $u$-morphismes. Si l'on prend des flèches $u$-cartésiennes $Y \to Z$, $Y' \to Z'$ on tire de (*) un diagramme commutatif de $i$-morphismes
    $$
    \raisebox{2em}{\xymatrix{
        P \ar[r] \ar[d] & Y' \ar[d] \\
        X \ar[r] & Y
    }}
    \quad \text{(**)}
    $$
    Et alors : pour que le diagramme (*) soit \emph{cartésien dans $F$}, il faut et il suffit que le diagramme (**) soit \emph{cartésien dans la fibre $F_i$} (cela vient de ce que les changements de base de $F$ sont exacts à gauche).
    En particulier :
    \item[(ii)] Pour tout objet $i$ de $I$, le fonction d'inclusion $F_i \to F$ commute aux produits fibrés.
\end{enumerate}
}

\vskip .3cm
{ 
Proposition {\bf 4.4.2.} --- \it
\begin{enumerate}
    \item[(i)] Soit $G$ un $\cU$-préfaisceau su $F$. $G$ est un faisceau si et seulement si la restriction $G_{/F_i}$ est représentable (i.e. un faisceau sur $F_i$) pour tout $i \in \Ob(I)$.
    \item[(ii)] Pour toute section cartésienne $S$ de $F$, le préfaisceau $p(S)$ (4.3) est un faisceau sur $F$.
\end{enumerate}
}
\vskip .3cm

((ii) découle de (i). (i) résulte immédiatement de la définition de $T$, compte tenu de la remarque (ii) de 4.4.1.)

\vskip .3cm
{ 
Proposition (Sous-catégories génératrices de $F$) {\bf 4.4.3.} --- \it Soit $S$ une sous-catégorie pleine de $F$. Pour que $S$ engendre le site $F$ (i.e. que tout objet de $F$ puisse être recouvert par des objets de $S$ au sens de la topologie $T$) il faut et il suffit que pour tout $i \in \Ob(I)$, la fibre $S_i$ engendre le topos $F_i$ (clair).
}
\vskip .3cm

{
Proposition {\bf 4.4.4.} --- \it
\begin{enumerate}
    \item[a)] De (3), on déduit que le site $F$ admet une petite famille topologiquement génératrice. Donc \emph{la catégorie $\widetilde{F}$ de $\cU$-faisceaux sur $F$ est un $\cU$-topos.}
    \item[b)] Soit $n: F \to \widetilde{F}$ le foncteur naturel, composé du foncteur naturel $F \to \widehat{F}$ et du foncteur faisceau associé. (On a une bijection fonctorielle
    $$
    G(X) \isommap \Hom(n(X), G)
    $$
    $X \in \Ob(F)$, $G \in  \Ob(\widetilde{F})$).
\end{enumerate}
}
\vskip .3cm

Le foncteur $n$ n'est pas pleinement fidèle en général : les objets vides des fibres de $F$ sont tous couverts par la famille vide, mais forment une sous-catégorie pleine de $F$ équivalente à $I$. Néanmoins :

\vskip .3cm
{
Lemme. --- \it Soit $F^*$ la sous-catégorie pleine de $F$ définie comme voici : pour tout $i \in \Ob(I)$, la fibre $(F^*)_i$ est la catégorie des objets non-vides de $F_i$ (autrement dit, $F^*$ est formée des objets $X$ tels que $n(X)$ soit non-vide dans $\widetilde{F}$). La restriction de $n$ à $F^*$ est pleinement fidèle.
}
\vskip .3cm

{\it Démonstration}. Soient $X$, $Y$ deux objets de $F^*$, de projections $i$ et $j$ respectivement, et $(f_{\alpha}: X_{\alpha} \to X)$ une famille épimorphique de $F_i$. Soit
$$
(g_{\alpha}: X_{\alpha} \to Y)
$$
une famille de morphismes de $F$, telle que pour tout couple d'indices $\alpha$, $\beta$ on ait le diagramme commutatif :
$$
    \raisebox{2em}{\xymatrix{
        X_\alpha \times_X X_\beta \ar[r] \ar[d] & X_\beta \ar[d]^{g_\beta} \\
        X_\alpha \ar[r]_{g_\alpha} & Y
    }}
    \quad \text{(*)}
$$

Pour tout indice $\alpha$, soit $u_{\alpha}: i \to j$ la projection du morphisme $g_{\alpha}: X_{\alpha} \to Y$. Les diagrammes (*) donnent des diagrammes commutatifs de $I$ :
$$
    \xymatrix{
        i \ar@{=}[r]  \ar@{=}[d] & i \ar[d]^{u_\beta}\\
        i \ar[r]_{u_\alpha} & j
    }
$$
donc les $u_{\alpha}$ sont égaux à un certain $u: i \to j$. Prenons maintenant une image inverse $Y' \to Y$ de $Y$ par $u$, et soient $(h_{\alpha})$ les $i$-morphismes $X_{\alpha} \to Y'$ déduits des $u$-morphismes $g_{\alpha}: X_{\alpha} \to Y$.

Soit $g: X \to Y$ un morphisme. Pour que les diagrammes
$$
    \raisebox{2em}{\xymatrix{
        X_\alpha \ar[d]_{f_\alpha} \ar[dr]^{g_\alpha} \\
        X \ar[r]_g & Y
    }}
    \quad \text{(**)}
$$
soient commutatifs, il faut et il suffit que
\begin{enumerate}
    \item[1)] $\Pi(g) = u$ et 
    \item[2)] si l'on désigne par $h: X \to Y'$ le $i$-morphisme déduit de $f$, les diagrammes
    $$
    \xymatrix{
        X_\alpha \ar[d]_{f_\alpha} \ar[dr]^{h_\alpha} \\
        X \ar[r]_h & Y'
    }
    $$
    de $F_i$ soient commutatifs. Donc il existe un $g: X \to Y$ et un seul qui rend commutatifs les diagrammes (**), c.q.f.d.
\end{enumerate}

\vskip .3cm
{
Proposition (Familles épimorphiques de $\widetilde{F}$) {\bf 4.4.5.} --- \it Une famille $(G_{\alpha} \to G)$ de morphismes de $\widetilde{F}$ est épimorphique si et seulement si, pour tout $i \in \Ob(I)$, la famille de morphismes de faisceaux sur $F_i$ :
$$
(G_{\alpha/F_i} \to G_{/F_i})
$$
est épimorphique.
}
\vskip .3cm

En effet, d'après 4.4.4, (a), ces deux propriétés sont équivalentes à la suivante : pour tout $i \in \Ob(I)$, tout $X \in \Ob(F_i)$ et tout $x \in G(X)$ il existe une famille épimorphique $(X_{\lambda} \to X)$ de $F_i$ telle que chacun des $x_{/X_{\lambda}}$ se relève l'un des $G_{\alpha}$. 

\vskip .3cm
{
Proposition {\bf 4.4.6.} --- \it
\begin{enumerate}
    \item[a)] Le foncteur
    $$
    p: L \to \widetilde{F}
    $$
    (4.3 et 4.4.2, (ii)) définit un morphisme de topos
    $$
    \widetilde{F} \to L
    $$
    \item[b)] Pour tout $i \in \Ob(I)$, il existe un morphisme de topos
    $$
    u_i: F_i \to \widetilde{F}
    $$
    tel que le foncteur composé
    $$
    \widetilde{F} \xlongrightarrow{u^{-1}_i} F_i \to \widetilde{F_i}
    $$
    soit isomorphe au foncteur ``restriction à $F_i$''.
\end{enumerate}
}
\vskip .3cm

Cela découle de 4.4.5, compte tenu de (4.1.1) pour (a) et de (4.4.2, (i)) pour (b).

\vskip .3cm
{
Proposition {\bf 4.4.7.} --- \it Soient $i \in \Ob(I)$, $C \in \Ob(F_i)$ et $n: F \to \widetilde{F}$ le foncteur naturel. Si $C$ est connexe dans $F_i$, $n(C)$ est connexe dans $\widetilde{F}$.
}
\vskip .3cm
{
Corollaire. --- \it Si les $F_i$ sont localement connexes, $\widetilde{F}$ est localement connexe.
}
\vskip .3cm

{\it Démonstration (de la proposition)}. Supposons $n(C)$ somme directe de $G$, $H \in \Ob(\widetilde{F})$. Il existe alors des objets non-vides $(D_\lambda)$ de $F_i$ et une famille épimorphique $D_\lambda \to C$ de $F_i$, telle que chaque $n(D_\lambda) \to n(C)$ se factorise par $G$ ou par $H$. Si $n(D_\lambda) \to n(C)$ se factorise par $G$ et $n(D_\mu) \to n(C)$ par $H$, le produit fibré
$$
n(D_\lambda) \times_{n(C)} n(D_\mu) \isom n(D_\lambda \times_C D_\mu) \quad \text{(cf 4.4.1, (ii))}
$$
est vide dans $\widetilde{F}$, ou encore, ce qui revient au même, $D_\lambda \times_C D_\mu$ est vide dans $F_i$. On conclut par (1.4) que $G$ ou $H$ est vide.

\vskip .3cm
{
Proposition {\bf 4.4.8.} --- \it Soient $u: i \to j$ une flèche de $I$, $M$ un objet de $F_j$, $C$ et $D$ des objets connexes et non-vides de $F_i$ et $F_j$ respectivement, et $f: C \to D$ un $u$-morphisme. Si $D$ trivialise $M$, alors l'application de composition avec $f$ : 
$$
\Hom_j(D, M) \to \Hom_u(C, M)
$$
est \emph{bijective}.
}
\vskip .3cm

{\it Démonstration}. Pour tout ensemble $H$, on tire de $f: C \to D$ (par somme directe) un $u$-morphisme $H_C \to H_D$ qui rend commutatif le diagramme
$$
     \xymatrix{
         H_C \ar[r] \ar[d] & H_D \ar[d]\\
         C\ar[r] & D
     }
$$
ce diagramme est même cartésien dans $F$ (par 4.4.1, (i)) ; et puisque $C$ et $D$ sont connexes et non-vides on a un diagramme commutatif d'applications :
$$
     \xymatrix@C=1em{
         & H \ar[dl] \ar[dr] \\
         \Hom_D(D,H_D) \ar[r] & \Hom_D(C,H_D) \ar[r]^\sim &\Hom_C(C,H_C)
     }
$$
Si on prend $H = \Hom_j(D, M)$, le $D$-isomorphisme $H_D \to D \times M$ (2.1.3) fournit un nouveau diagramme commutatif:
$$
     \xymatrix{
         \Hom_D(D,H_D) \ar[r]^{\sim} \ar[d]& \Hom_D(C,H_D) \ar[d]\\
         H = \Hom_j(D,M) \ar[r] & \Hom_u(C,M)
     }
$$
d'où notre proposition.

\vskip .3cm
{
Proposition (Image essentielle du foncteur $p$ (4.3)) {\bf 4.4.9.} --- \it Soit $G$ un faisceau sur $F$. Pour tout $i \in \Ob(I)$ représentons la restriction $G_{/F_i}$ par un objet $X_i$ de $F_i$, et soit $\epsilon_i \in G(X_i)$ la section qui définit l'isomorphisme $X_i \to G_{/F_i}$. Les propriétés suivantes sont équivalentes :
\begin{enumerate}
    \item[(i)] Il existe une section cartésienne $Y: I \to F$ et un isomorphisme $p(Y) \isommap G$ ; 
    \item[(ii)] Pour toute flèche $u: i \to j$ de $I$, il existe un morphisme $u$-cartésien $X_i \xlongrightarrow{m} X_j$ qui rend commutatif le diagramme :
    $$
    \raisebox{2em}{\xymatrix@C=1em{
        n(X_i) \ar[dr]_{\epsilon_i} \ar[rr]^{n(m)} && \ar[dl]^{\epsilon_j} n(X_j)\\
        & G
    }}
     \quad \text{(*)} $$
    (où $n$ désigne le foncteur   naturel $F \to \widetilde{F}$).
\end{enumerate}
}
\vskip .3cm

{\it Démonstration}. (i) $\Rightarrow$ (ii) Pour tout $i \in \Ob(I)$, désignons par $\alpha_i$ la section $1_{Y_i}$ de $p(Y)$ au-dessus de $Y(i)$ (4.3). Il existe un isomorphisme et un seul $Y(i) \to X_i$ qui rend commutatif le diagramme
 $$
     \xymatrix{
         n(Y(i)) \ar[r]^-\sim \ar[d]_{\alpha_i} & n(X_i) \ar[d]^{\epsilon_i}\\
         p(Y) \ar[r]^\sim & G
     }
 $$
 Or, pour toute flèche $u: i\to j$ de $I$, le diagramme
 $$
     \xymatrix@C=1em{
         n(Y(i)) \ar[dr]_{\alpha_i} \ar[rr] && n(Y(j)) \ar[dl]^{\alpha_j}\\
         & p(Y)
     }
 $$
 est commutatif, d'où (ii).
 
 (ii) $\Rightarrow$ (i) Soit $u: i \to j$ une flèche de $I$. Prouvons qu'il n'existe qu'une flèche $u$-cartésienne $m: X_i \to X_j$ qui rend commutatif le diagramme (*) : soit $m'$ une seconde flèche qui jouit de la même propriété. Il existe un $i$-automorphisme de $X_i$ et un seul, soit $\alpha$, tel que $m' = m \circ \alpha$. On a donc :
 $$
 \epsilon_i = \epsilon_j \circ n(m') = \epsilon_j \circ n(m) \circ n(\alpha) = \epsilon_i \circ n(\alpha)
 $$
 cela veut dire que $1_{X_i}$ et $\alpha$ ont même image par l'isomorphisme $X_i \isommap G_{/F_i}$\footnote{Plus précisément, par la bijection composée $\Hom_i(X_i, X_i) \isommap G(X_i) \isommap \Hom_{\widetilde{F}}(n(X_i), G)$}, d'où $m = m'$.
 
 Nous pouvons donc définir une section cartésienne $X: I \to F$ comme voici :
 \begin{enumerate}
     \item[a)] $X(i) = X_i$ pour tout $i \in \Ob(I)$
     \item[b)] pour toute flèche $i \xlongrightarrow{u} j$ de $I$, la flèche $u$-cartésienne $X(u): X_i \to X_j$ est déterminée par le diagramme commutatif (*) (en effet, l'unicité que l'on vient de prouver nous garantit que $X(v \circ u) = X(v) \circ X(u)$ pour tout couple de flèches $i \xlongrightarrow{u} j \xlongrightarrow{v} k$ de $I$). Et il est clair que $p(X)$ est isomorphe à $G$.
 \end{enumerate}

\vskip .3cm
{\bf 4.5. Théorème de Van Kampen}

Comme dans (2.4) et (3) nous supposons que l'univers $\cU$ admet un élément de cardinal infini.

\vskip .3cm
{
Théorème. --- \it Si les fibres de $F$ sont localement connexes, les topos $\widetilde{F}$ et $L$ sont localement connexes, et le morphisme de topos
$$
\widetilde{F} \to L
$$
(4.4.6, a) défini par le foncteur pleinement fidèle $p: L \to \widetilde{F}$ (4.3) fournit une équivalence de topos
$$
\SLC(\widetilde{F}) \to \SLC(L) \quad (2.4.12), \quad (3.4)
$$
}
\vskip .3cm

{\it Remarque}. cela   revient à dire que le foncteur $p: L \to \widetilde{F}$ induit par restriction une équivalence entre catégories d'objets localement constants. Ainsi la démonstration du théorème se ramène aux points 4.5.1 et 4.5.2 qui suivent :

\vskip .3cm
{\bf 4.5.1.} Soit $G$ un faisceau sur $F$. Si $G$ est localement constant dans $\widetilde{F}$, il existe une section cartésienne $X: I \to F$ telle que $p(X)$ soit isomorphe à $G$.

Nous utiliserons le critère 4.4.9, dont nous reprenons les notations.

Soit $u: i \to j$ une flèche de $I$.

Soit $K$ l'ensemble des couples
$$
(C, D) \in \Ob(F_i) \times \Ob(F_j)
$$
qui remplissent les conditions suivantes :
\begin{enumerate}
    \item[(i)] $C$ et $D$ sont connexes et non-vides dans $F_i$ et $F_j$ respectivement
    \item[(ii)] $\Hom_u(C,D) \neq \emptyset$
    \item[(iii)] $n(D)$ trivialise $G$
\end{enumerate}
Soit enfin $S$ la sous-catégorie pleine de $F_i$ formée des objets $C$ qui remplissant la condition suivante : il existe un objet $D$ de $F_j$ tel que $(C, D) \in K$.
\begin{enumerate}
    \item[a)] \emph{$S$ engendre $F_i$ :}
    
    En effet (4.4.3) les objets connexes $Z$ de $F_j$ tels que $n(Z)$ trivialise $G$ engendrent $F_j$. Donc leurs images inverses par $u$ \emph{recouvrent} $F_i$ ; d'où notre assertion.
    \item[b)] \emph{Soit $C$ un objet de $S$. Pour chaque $x \in G(C)$, il existe un $u$-morphisme $C \to X_j$ et un seul qui rend commutatif le diagramme :}
    $$
    \xymatrix{
        & n(X_j) \ar[d]^{\epsilon_j}\\
        n(C) \ar[ur] \ar[r]_x & G
    }
    $$
    {\it Démonstration}. Soient $D$ un objet de $F_j$ tel que $(C, D) \in K$ et $f: C \to D$ un $u$-morphisme. On en tire un diagramme commutatif~:
$$
    \xymatrix{
        \Hom_j(D,X_j) \ar[r] \ar[d] & G(D) \ar[d]\\
        \Hom_u(C,X_j) \ar[r] & G(C)
    }  
$$
(où les flèches verticales se déduisent de $C \to D$ et les flèches horizontales de $n(X_j) \to G$).

Il s'agit de prouver que la flèche du bas est bijective ; or :
\begin{enumerate}
    \item[(i)] La flèche du haut est bijective par hypothèse ($X_j$ représente $G_{/F_j}$)
    \item[(ii)] La flèche $\Hom_j(D, X_j) \to \Hom_u(C, L_j)$ est bijective par (4.4.8) : en effet, la restriction $n(D)_{/F_j}$ trivialise $X_j$ d'après (4.4.6, b), et elle admet une section au-dessus de $D$.
    \item[(iii)] La flèche $G(D) \to G(C)$ est bijective parce que $n(C)$ et $n(D)$ sont connexes non-vides dans $\widetilde{F}$ (4.4.7) et que $n(D)$ trivialise $G$.
\end{enumerate}
    \item[c)] Pour tout $C \in \Ob(S)$ et tout $f \in \Hom_i(C, X_i)$, soit $\alpha_C(f)$ le $u$-morphisme $C \to X_j$ qui rend commutatif le diagramme
    $$
    \xymatrix@C=1em{
        n(C) \ar[d] \ar[drr]\\
        n(X_i) \ar[dr] && n(X_j) \ar[dl]\\
        & G
    }
    $$
    L'application 
    $$
    \alpha_C: \Hom_i(C, X_i) \to \Hom_u(C, X_j)
    $$
    est nécessairement fonctorielle en $C$ ; elle est bijective puisque l'application
    $$
    \Hom_i(C, X_i) \to G(C)
    $$
    fournie par $X_i \to G_{/F_i}$ l'est. Puisque $S$ engendre $F_i$, les bijections $\alpha_C$ proviennent d'une flèche $u$-cartésienne
    $$
    X_i \to X_j
    $$
    et le diagramme
    $$
    \xymatrix@C=1em{
        n(X_i) \ar[rr] \ar[dr] && n(X_j) \ar[dl]\\
        & G
    }
    $$
    est commutatif.
\end{enumerate}

\vskip .3cm
{
Proposition {\bf 4.5.2.} --- \it Soit $X: I \to F$ une section cartésienne. Si $p(X)$ est localement constant dans $\widetilde{F}$, alors $X$ est un objet localement constant de $L$.
}
\vskip .3cm

{\it Démonstration}. Recouvrons l'objet final de $\widetilde{F}$ par des objets localement constants $(H_{\alpha})$ qui trivialisent $p(X)$ (par exemple des objets galoisiens 2.4.6). Si on prend des sections cartésiennes $(M_{\alpha})$ de $F$ telles que les $p(M_{\alpha})$ soient isomorphes aux $H_{\alpha}$, les $M_{\alpha}$ trivialisent $X$ puisque $p$ est pleinement fidèle ; or les $M_{\alpha}$ recouvrent $e_L$ par 4.4.5.

\vskip .3cm
{\bf 4.6. Étude de $\SLC(\widetilde{F}) \isom \SLC(L)$} 
(Nous supposons donc que l'univers $\cU$ admet un élément de cardinal infini, et que les fibres de $F$ sont localement connexes)

Rappelons certaines notations :
\begin{enumerate}
    \item[1)] $n: F \to \widetilde{F}$ le foncteur naturel (4.4.4)
    \item[2)] $p: L \to \widetilde{F}$ le foncteur pleinement fidèle qui définit le morphisme de topos $\widetilde{F} \to L$ (4.4.6).
\end{enumerate}

\vskip .3cm
{
Proposition {\bf 4.6.1.} --- \it Soit $H$ un objet localement constant de $\widetilde{F}$. Il existe une sous-catégorie pleine $C$ de $F$ qui remplit les conditions suivantes :
\begin{enumerate}
    \item[1)] Pour tout $X \in \Ob(C)$, $n(X)$ trivialise $H$.
    \item[2)] Pour tout $i \in \Ob(I)$, les objets de $C_i$ sont galoisiens dans $F_i$ et recouvrent $e_{F_i}$. En outre, deux objets de $C_i$ situés au-dessus de la même composante connexe de $F_i$ sont toujours isomorphes (cette dernière condition veut dire que $C_i$ est un groupoïde (2.3.4)).
    \item[3)] Pour tout flèche $u: i \to j$ de $I$ et tout objet $Y$ de $C_i$, il existe un objet $Z$ de $C_j$ et un $u$-morphisme $Y \to Z$.
\end{enumerate}
}
\vskip .3cm

{\it Démonstration}. Soient $(U_\alpha)_{\alpha \in A}$ les composantes connexes de $e_{\widetilde{F}}$. Au-dessus chaque $U_\alpha$, prenons un objet galoisien $G_\alpha$ qui trivialise $H$ (2.4.6).

Il existe (4.5) des sections cartésiennes $S_\alpha$ de $F$ et des isomorphismes $p(S_\alpha) \isom F$. Nous définissons la catégorie $C$ comme voici : pour tout $i \in \Ob(I)$, la fibre $C_i$ est la sous-catégorie pleine de $C_i$ formée des composantes connexes des $S_\alpha(i)$.

Vérification de (1) : pour tout $X \in \Ob(C)$, il existe un $\alpha \in A$ et une section de $G_\alpha$ au dessus de $X$, i.e. un morphisme $n(X) \to G_\alpha$. Donc $n(X)$ trivialise $H$.

Vérification de (2) : nous appliquons (4.4.6, b) dont nous reprenons les notations. Soit $i \in \Ob(I)$. Les $u^{-1}_i(G_\alpha)$ sont isomorphes aux $S_\alpha(i)$. Par conséquent a) pour tout $\alpha \in A$, $S_\alpha(i)$ admet une structure de pseudo-torseur sous le groupe constant $\Aut(G_\alpha)_{F_i}$ et b) les $S_\alpha(i)$ recouvrent $e_{F_i}$ et $S_\alpha(i) \times S_\beta(i)$ est vide pour $\alpha \neq \beta$. Le point (2) découle donc de 2.4.5., (ii).

Quant au point (3), il est évident.

\vskip .3cm
{
Proposition {\bf 4.6.2.} --- \it Soit $C$ une sous-catégorie pleine de $F$ qui remplit la condition (2) de (4.6.1). Pour tout faisceau $H$ sur $F$, les propriétés suivantes sont équivalentes :
\begin{enumerate}
    \item[(i)] Pour tout $X \in \Ob(C)$, $n(X)$ trivialise $H$.
    \item[(ii)] Il existe une section cartésienne $S$ de $F$ telle que $p(S)$ soit isomorphe à $H$ et que, pour tout $i \in \Ob(I)$, les objets de $C_i$ trivialisent $S(i)$.
\end{enumerate}
}
\vskip .3cm
{
Corollaire. --- \it Soit $S$ une section cartésienne de $F$. Pour que $S$ soit un objet localement constant de $L$, il faut et il suffit que pour tout $i \in \Ob(I)$, $S(i)$ soit un objet localement constant de $F_i$ (cf. 2.4.6 et 4.6.1).
}
\vskip .3cm

{\it Démonstration de la proposition}. (i) $\Rightarrow$ (ii) Par définition de la topologie de $F$, les $n(X)$, $X \in \Ob(C)$ recouvrent le faisceau final. Donc $H$ est localement constant, et il existe une section cartésienne $S$ de $F$ telle que $p(S)$ soit isomorphe à $H$. Soient $i \in \Ob(I)$ et $X$ un objet de $C_i$ ; la restriction $n(X)_{/F_i}$ admet une section au-dessus de $X$, et elle trivialise $H_{/F_i}$ d'après (4.4.6, b) ; donc $X$ trivialise $S(i)$ qui représente $H_{/F_i}$, c.q.f.d. 

(ii) $\Rightarrow$ (i) Soit $K$ la sous-catégorie pleine de $F$ définie comme voici : pour tout $i \in \Ob(I)$, la fibre $K_i$ est formée des objets connexes non-vides de $F_i$ qui sont dans le crible engendré par les objets de $C_i$. Le foncteur $n: F \to \widetilde{F}$ identifie $K$ à une sous-catégorie génératrice de $\widetilde{F}$ formée d'objets connexes non-vides (d'après le lemme 4.4.4, b et 4.4.7). Il suffit donc, d'après (2.2.1), de voir que $H$ est un préfaisceau localement constant sur $K$ ; ce qui découle de (4.4.8), compte tenu de l'isomorphisme $p(S) \isom H$. 

\vskip .3cm
{
Proposition {\bf 4.6.3.} --- \it Soit $C$ une sous-catégorie pleine de $F$ qui remplit les conditions (2) et (3) de (4.6.1).
\begin{enumerate}
    \item[(i)] Pour tout $i \in \Ob(I)$, soit $\LC(F_i, C_i)$ la sous-catégorie pleine de $F_i$ formée des objets qui sont trivialisés par les objets de $C_i$. $\LC(F_i, C_i)$ est un $\cU$-topos objet de $\cG$ (cf. 3.2), et l'inclusion $\LC(F_i, C_i) \to F_i$ définit un morphisme de topos
    $$
    F_i \to \LC(F_i, C_i) \quad \text{(3.3.2)}
    $$
    Le topos $\LC(F_i, C_i)$ est engendré par $C_i$. Comme $C_i$ est un groupoïde, $\LC(F_i, C_i)$ est équivalent à $\widehat{C}_i$.
    \item[(ii)] Soit $F_C$ la sous-catégorie pleine de $F$ qui a pour fibres les $\LC(F_i, C_i)$. $F_C \to I$ est une sous-catégorie fibrée en $\cU$-topos objets de $\cG$, et l'inclusion $F_C \to F$ définit un morphisme de catégories fibrées en topos 
    $$
    F \to F_C
    $$
    (preuve: d'après (i), il suffit de voir que $F_C$ est une sous-catégorie fibrée de $F$. Soient donc $u: i \to j$ une flèche de $I$ et $L \to M$ une flèche $u$-cartésienne de $F$. D'après la condition (3), si les objets de $C_j$ trivialisent $M$, les objets de $C_i$ trivialisent $L$).
    \item[(iii)] L'inclusion $F_C \to F$ identifie donc les sections cartésiennes de $F_C$ à des sections cartésiennes de $F$. D'après (4.6.2), le foncteur
    $$
    p: L \to \widetilde{F}
    $$
    fournit une \emph{équivalence} de catégories
    $$
    \Cart_I(I, F_C) \to \LC(\widetilde{F}, C)
    $$
    entre les sections cartésiennes de $F_C$ et les $\cU$-faisceaux sur $F$ qui sont trivialisés par les $n(X)$, $X \in \Ob(C)$.
    
    Ainsi \emph{le topos $\LC(\widetilde{F}, C)$ (3.3.2) est limite inductive des topos $\LC(F_i, C_i)$} (4.1.1) ; et l'inclusion $\Cart_I(I, F_C) \to \SLC(L)$ définit un morphisme de topos 
    $$
    \SLC(L) \to \Cart_I(I, F_C)
    $$
\end{enumerate}
}

\vskip .3cm
{\bf 4.6.4.} (Avec les hypothèses et les notations de 4.6.3)

On peut aussi définir sur $F_C$ la topologie (4.4.1). Le foncteur naturel
$$
n_C: F_C \to \widetilde{F}_C
$$
identifie alors la sous-catégorie pleine $C$ de $F_C$ à une sous-catégorie génératrice du topos $\widetilde{F}_C$ formée d'objets connexes et non-vides ((4.6.3, (i)), (4.4.3), (4.4.4 (b)) et (4.4.7)).

Désignons par :
$$
p_C: \Cart_I(I, F_C) \to \widetilde{F}_C
$$
le foncteur (4.3).

\vskip .3cm
{
Lemme. --- \it Soit $G$ un faisceau sur $F_C$. Les propriétés suivantes de $G$ sont équivalentes :
\begin{enumerate}
    \item[a)] $G$ est localement constant dans $\widetilde{F}_C$
    \item[b)] La restriction $G_{/C}$ est un préfaisceau localement constant
    \item[c)] Il existe une section cartésienne $S$ de $F_C$ et un isomorphisme $p_C(S) \isom G$.
\end{enumerate}
}
\vskip .3cm

{\it Démonstration}. On prouve $a \Rightarrow c \Rightarrow b \Rightarrow a$ :

Si $G$ est localement constant, il existe (d'après le théorème 4.5 appliqué à $F_C$) une section cartésienne $S$ de $F_C$ et un isomorphisme $p_C(S) \isom G$. Mais alors, pour tout $i \in \Ob(I)$, $S(i)$ est trivialisé par les objets de $C_i$ ; donc (par (4.6.2) appliqué à $F_C$ et (2.2.1)) $G_{/C}$ est un préfaisceau localement constant. L'implication $b \Rightarrow a$ découle de 2.2.1.

\vskip .3cm
{
Corollaire (de 4.6.4) {\bf 4.6.5.} --- \it
Le foncteur ``restriction à $C$'' donne une équivalence de catégories
$$
\LC(\widetilde{F}, C) \to \LC(\widehat{C})
$$
où $\LC(\widehat{C})$ désigne la catégorie des $\cU$-préfaisceaux localement constants sur $C$.
}
\vskip .3cm

(Appliquer (4.6.3, iii) et l'équivalence $b \Leftrightarrow c$ du lemme (4.6.4))

{\it Commentaire} : 
\begin{enumerate}
    \item[1)] Les préfaisceaux localement constants sur $C$ s'identifient aux préfaisceaux sur les groupoïde fondamental $\Pi(C)$ de la catégorie $C$, obtenu par calcul des fractions (\cite{1}, chap. I, n$^\circ 1$.5.3). Rappelons que $\Pi(C)$ est caractérisé à équivalence près par la propriété universelle suivante : il existe un foncteur $C \to \Pi(C)$ qui donne pour tout groupoïde $\Gamma$ une équivalence de catégories :
    $$
    \Fonct(\Pi(C), \Gamma) \to \Fonct(C, \Gamma)
    $$
    \item[2)] Le groupoïde $\Pi(C)$ est la 2-limite inductive des groupoïdes $C_i$ au sens suivant :
    \begin{enumerate}
        \item[a)] Le foncteur $\Pi_{/C}: C \to I$ est cofibrant, i.e. le foncteur $C^\circ \to I^\circ$ qui s'en déduit est fibrant (condition (3) de (4.6.1) et (4.4.8))
        \item[b)] $\Pi(C)$ représente le 2-foncteur
        $$
        \Gamma \to \Fonct_I(C, I \times \Gamma)
        $$
        des groupoïdes dans les catégories.
    \end{enumerate}
    \item[3)] On peut exprimer cela autrement : d'après le corollaire (4.6.2), la sous-catégorie $\cG$ de $\Top$ définie en (3.2) est stable par les 2-limites inductives (les objets de $\cG$ sont d'ailleurs les limites inductives de topos ponctuels). Or le foncteur:
    $$
    \Point: \cG \to \Grpd
    $$
    commute aux 2-limites inductives (puisque c'est une équivalence) ; et dans le cas présent, les $C_i$ sont équivalents aux groupoïdes $\Point(\LC(F_i, C_i))$ (3.2.4 et 4.6.3, (i)). On peut donc se reporter à 4.6.3, (iii).
\end{enumerate}

\vskip .3cm
{
Proposition {\bf 4.6.6.} --- \it Soit $J(F)$ l'ensemble des sous-catégories pleines de $F$ qui remplissent les conditions (2) et (3) de 4.6.1. On définit une relation de préordre sur $J(F)$ comme voici : $C < D$ si pour tout $i \in \Ob(I)$, les objets de $C_i$ trivialisent les objets de $D_i$.
\begin{enumerate}
    \item[(i)] L'ensemble préordonné $J(F)$ est filtrant à gauche
    \item[(ii)] L'ordonné associé est$\cU$-petit.
\end{enumerate}
}
\vskip .3cm

{\it Démonstration}. Le point (ii) résulte de 2.4.9.

Prouvons (i). Soient $C$ et $C'$ deux éléments de $J(F)$. On définit un minorant $D$ de $\{C, C'\}$ comme voici : étant donné $i \in \Ob(I)$, un objet $Z$ de $F_i$ est dans $D_i$ s'il existe un objet $Y$ de $C_i$ et un objet $Y'$ de $C_i'$ tels que $Z$ soit isomorphe à une composante connexe de $Y \times Y'$ (cf. 2.4.5, (ii)). 

\vskip .3cm
{\bf 4.6.7.} D'après (4.6.6), (4.6.1) et (4.6.2), il existe un petit ensemble ordonné filtrant à gauche $A$ et une famille croissante $(C^\alpha)_{\alpha \in A}$ d'éléments de $J(F)$ qui remplit la condition suivante : pour tout objet localement constant $S$ de $L$, il existe un indice $\alpha$ tel que pour tout $i \in \Ob(I)$, les objets de $C^\alpha_i$ trivialisent $S(i)$ dan $F_i$. Cela veut dire que la catégorie des objets localement constants de $L$ est réunion des sous-catégories essentiellement pleines
$$
T_{\alpha} = \Cart_I(I, F_{C^ \alpha})
$$
définies en 4.6.3, (ii) et (iii). Puisque les $T_{\alpha}$ sont des $\cU$-topos objets de $\cG$ et que les inclusions $T_{\alpha} \to \SLC(L)$ définissent des morphismes de topos $\SLC(L) \to T_{\alpha}$, on a, suivant (3.3.3), la ``formule''
$$
\SLC(\varinjlim_I F_i) = \varprojlim_A (\varinjlim_I \LC(F_i, C^\alpha_i))
$$
où les $\lim$ sont des 2-limites dans la 2-catégorie des $\cU$-topos.


% End

% Begin


%%%%%%%%%%%%%%%%%%%%%%%%%%%%%%%%%%%%%%%%%%%%%%%%%%%%%%%%%%%%%%%
\chapter*{\S \space V. --- COMPLÉMENTS}
\addcontentsline{toc}{chapter}{V. Compléments}
\section*{}

{\bf 5.1. Groupe fondamental d'un topos localement connexe en un point}
\vskip .3cm
{
Définition {\bf 5.1.1.} --- \it Soient $T$ un topos \emph{localement galoisien} et$p$ un point de $T$. Nous appellerons \emph{groupe fondamental de $T$ en $p$} le groupe $\Pi_1 = \Pi_1(T, p)$ des automorphismes du foncteur fibre $p^{-1}$. On a donc pour tout objet $X$ de $T$ une opération à gauche naturelle de $\Pi_1(T, p)$ sur la fibre $p^{-1}(X)$ ; c'est-à-dire un foncteur
$$
f: T \to \Ens_{\Pi_1}
$$
de $T$ dans la catégorie des $\Pi_1$-ensembles à gauche.
}

\vskip .3cm
{
Proposition {\bf 5.1.2.} --- \it Pour tout objet localement constant $L$ de $T$, soit $V_L$ l'ensemble des $\alpha \in \Pi_1$ qui laissent fixe chaque point de $p^{-1}(L)$. Les ensembles $V_L$ forment un système fondamental de voisinages de $1$ pour une topologie de groupe sur $\Pi_1$. Pour tout objet $X$ de $T$, l'opération de $\Pi_1$ sur $p^{-1}(X)$ est alors continue pour la topologie discrète de $p^{-1}(X)$ ; d'où un nouveau foncteur
$$
\overline{f}: T \to \Dis_{\Pi_1} 
$$
à valeurs dans la catégorie des $\Pi_1$-espaces discrets. 
}

\vskip .3cm
{\bf 5.1.3.} Soit $I$ la catégorie des voisinages galoisiens de $p$ : les objets de $I$ sont les couples $(Y, y)$ formés d'un objet galoisien $Y$ de $T$ et d'un $y \in p^{-1}(Y)$ ; et les morphismes $(Y, y) \to (Z, z)$ sont les $Y \to Z$ qui transforment $y$ en $z$. $I$ est en fait un ensemble préordonné filtrant. Pour toute flèche $u: (Y, y) \to (Z, z)$ de $I$, on définit un morphisme de groupes surjectif $\Aut(Y) \to \Aut(Z)$ en associant à l'automorphisme $a$ de $Y$ l'automorphisme $b$ de $Z$ tel que $b(z) = u(a(y))$ ; d'où un système projectif de groupes discrets
$$
\Aut(Y)_{(Y, y) \in I}
$$

Le groupe topologique $\Pi_1$ s'identifie à la limite projective de ce système si on fait correspondre à chaque $\alpha \in \Pi_1$ la famille $(a_{(Y, y)})$ déterminée par les relations
$$
a_{(Y, y)}(\alpha\cdot y) = y
$$

\vskip .3cm
{
Proposition {\bf 5.1.4.} --- \it Supposons maintenant le topos $T$ \emph{connexe}. Les propositions suivantes sont alors équivalentes :
\begin{enumerate}
    \item[(a)] Le foncteur $\overline{f}: T \to \Dis_{\Pi_1}$ est une \emph{équivalence de catégories} 
    \item[(b)] Pour tout voisinage galoisien $(Y, y)$ de $p$, la projection
    $$
    \pr_{(Y, y)}: \Pi_1 \to \Aut(Y)
    $$
    est \emph{surjective}.
    \item[(b')] Pour tout objet connexe $M$ de $T$, $\Pi_1$ opère transitivement sur $p^{-1}(M)$.
\end{enumerate}
}
\vskip .3cm
Ces propositions sont vérifiées dans les deux cas suivants :
\begin{enumerate}
    \item[(i)] Tout objet galoisien de $T$ est fini (cf. 5.2.)
    \item[(ii)] $T$ admet une famille génératrice dénombrable.
\end{enumerate}
Notons qu'un topos localement galoisien qui remplit la condition (i) ou la condition (ii) admet toujours un point.

\vskip .3cm
{
Définition {\bf 5.1.5.} --- \it Soit maintenant $E$ un topos localement connexe. A chaque point $p$ de $E$, le morphisme de topos $E \to \SLC(E)$ (2.4.1.) fait correspondre un point $\overline{p}$ de $\SLC(E)$.

Le foncteur fibre $p^{-1}$ n'est autre que la restriction de $p^{-1}$ à la sous-catégorie $\SLC(E)$ de $E$. On peut appeler groupe fondamental de $E$ en $p$ le groupe fondamental en $\overline{p}$ du topos localement galoisien $\SLC(E)$.
}

Les propriétés (i) et (ii) de 5.1.4., pour le topos $\SLC(E)$, reviennent aux propriétés suivantes de $E$ :
{\it
\begin{enumerate}
    \item[(i')] Tout objet galoisien de $E$ est fini.
    \item[(ii')] Il existe une \emph{suite} $(R_n)$ de cribles couvrants $e_E$, telle que chaque objet localement constant de $E$ puisse être trivialisé par un $R_n$ (cf. 3.3.2. et 3.2.5.)
\end{enumerate}
}

{\bf 5.2. Groupoïde fondamental profini}

\vskip .3cm
{
Proposition {\bf 5.2.1.} --- \it Disons qu'un objet localement constant $L$ d'un topos $T$ est \emph{fini} s'il existe un recouvrement $(U_{\alpha})$ de $e_T$ par des objets de $T$ et des $U_{\alpha}$-isomorphismes
$$
U_{\alpha} \times L \isom I^{\alpha}_{U_{\alpha}}
$$
où les $I^{\alpha}$ sont des ensembles finis.

On prouve sans peine les propositions suivantes : 
\begin{enumerate}
    \item[1)] Soient $L$ un objet l.c.f. de $T$ et $R$ une relation d'équivalence sur $L$. Si $R$ est un objet l.c.f., le quotient l'est aussi.
    \item[2)] Toute limite projective finie d'objets l.c.f. est l.c.f.
\end{enumerate} 
Et, \emph{si $T$ est somme directe de topos connexes} :
\begin{enumerate}
    \item[3)] Tout objet $L$ de $T$ qui est l.c.f. est somme directe d'objets l.c.f. connexes ; et tout sous-objet de $L$ qui est l.c.f. est somme directe de composantes connexes de $L$.
    \item[4)] Soit $L$ un objet l.c.f. de $T$. Au-dessus de chaque composante connexe de $T$, il y a un objet galoisien fini qui trivialise $L$.
\end{enumerate}
}

\vskip .3cm
{\bf 5.2.2.} Nous supposons le topos $T$ somme directe de topos connexes, et que l'univers de référence admet un élément infini.

Soit $\SLCF(T)$ la sous-catégorie pleine de $T$ formée des sommes directes d'objets l.c.f. Les propositions 1 à 3 ci-dessus montrent que la catégorie $K$ des objets l.c.f. de $T$ vérifie les hypothèses du lemme 2.4.2.. Or, on prouve aisément que cette catégorie est petite à équivalence près ; donc $\SLCF(T)$ est un topos et l'inclusion $\SLCF(T) \to T$ définit un morphisme de topos en sens inverse. Enfin, d'après la proposition 4 ci-dessus et le lemme 2.4.10., $\SLCF(T)$ est un topos localement galoisien.

\vskip .3cm
{
Proposition {\bf 5.2.3.} --- \it Disons qu'un topos localement galoisien est \emph{profini} s'il est engendré par ses objets galoisiens finis. Le topos localement galoisien $\SLCF(T)$ est profini, et le morphisme $T \to \SLCF(T)$ fournit pour tout topos localement galoisiens profini $P$ une équivalence de catégories
$$
\Homtop(\SLCF(T), P) \to \Homtop(T, P)
$$
}



% End

%%%%%%%%%%%%%%%%%%%%%%%%%%%%%%%%%%%%%%%%%%%%%%%%%%%%%%%%%%%%%%%
\chapter*{APPENDICE: CATÉGORIES FIBRÉES EN TOPOS}\thispagestyle{empty}
\label{sec:app}
\section*{}

Soit $I$ une catégorie $\cU$-petite.

\vskip .3cm
{\bf A.1.} J'appelle catégorie fibrée en $\cU$-topos au dessus de $I$ toute catégorie fibrée $F \xlongrightarrow{\Pi} I$ qui vérifie les axiomes suivantes : 
\begin{enumerate}
    \item[1)] Pour tout $i \in \Ob(I)$, la fibre $F_i$ est un $\cU$-topos. 
    \item[2)] Pour toute flèche $u: i \to j$ de $I$, le foncteur changement de base $F_j \to F_i$ définit un morphisme de topos $F_i \to F_j$.
\end{enumerate}

\vskip .3cm
{\bf A.2.} Définissons maintenant la $2$-catégorie $\Fibtop(I)$ des catégories fibrées en $\cU$-topos au-dessus de $I$ : étant données deux catégories fibrées en topos $F$, $G$ au-dessus de $I$, nous prenons comme catégorie des morphismes de $F$ dan $G$
$$
\Cartop_I(F, G)
$$
la sous-catégorie pleine de 
$$
\Cart_I(G, F)^{\circ}
$$
(catégorie opposée de la catégorie des $I$-foncteurs cartésiens $G \to F$) définie comme voici~: un $I$-foncteur cartésien $\phi: G \to F$ définit un morphisme de catégories fibrées en topos $F \to G$ si pour tout $i \in \Ob(I)$ le foncteur $G_i \to F_i$ déduit de $\phi$ par restriction définit un morphisme de topos $F_i \to G_i$.

\vskip .3cm
{\bf A.3.} La $2$-catégorie $\Fibtop(I)$ est équivalente à la $2$-catégorie des $2$-foncteurs de $I$ dans la catégorie des $\cU$-topos: les catégories fibrées de la forme $I \times E$ ($E$ un $\cU$-topos) correspondant aux $2$-foncteurs constants ; d'où une définition des $2$-limites inductives et projectives de topos~:

\vskip .3cm
{\bf A.4.} Soit $F$ une catégorie fibrée en $\cU$-topos au-dessus de $I$. Nous appellerons $2$-limite inductive de $F$ le $2$-foncteur covariant
\begin{align*}
 \cU\text{-topos} \to \cU\text{-catégories}\\
 E \mapsto \Cartop_I(F, I \times E)
\end{align*}
et $2$-limite projective de $F$ le $2$-foncteur contravariant
$$
E \to \Cartop_I(I \times E, F)
$$
La $2$-limite inductive (resp. projective) de $F$ se représente donc, quand c'est possible, par un $\cU$-topos $L$ muni d'un morphisme de catégories fibrées en topos
$$
F \to I \times L
$$
$$
\text{(resp.}~I \times L \to F).
$$


%%%%%%%%%%%%%%%%%%%%%%%%%%%%%%%%%%%%%%%%%

\newpage
\mbox{}
\thispagestyle{empty}

\end{document}


%End
